\documentclass[11pt,a4paper]{article}
\usepackage[T1]{fontenc}
\usepackage[utf8]{inputenc}
\usepackage{lmodern}
\usepackage[margin=27mm]{geometry}
\usepackage{amsmath,amssymb,amsthm,mathtools}
\usepackage{mathrsfs}
\usepackage{microtype}
\usepackage{booktabs,longtable,tabularx,array}
\usepackage{xcolor}
\usepackage{enumitem}
\usepackage{fancyhdr}
\usepackage{xurl}
\usepackage{hyperref}
\hypersetup{colorlinks=true,linkcolor=blue!45!black,citecolor=blue!45!black,
 urlcolor=blue!45!black,pdftitle={Exact Polylogarithm Parity, Harmonic-Tail Reconstruction, and Higher Polygamma Moments},
 pdfauthor={Research continuation for the ProveIt project}}
\setlength{\emergencystretch}{2em}
\setlength{\parskip}{3pt}
\setlength{\parindent}{1.1em}
\setlist{itemsep=3pt,topsep=5pt}
\numberwithin{equation}{section}
\newtheorem{theorem}{Theorem}[section]
\newtheorem{lemma}[theorem]{Lemma}
\newtheorem{corollary}[theorem]{Corollary}
\newtheorem{proposition}[theorem]{Proposition}
\newtheorem{conjecture}[theorem]{Conjecture}
\theoremstyle{definition}
\newtheorem{definition}[theorem]{Definition}
\theoremstyle{remark}
\newtheorem{remark}[theorem]{Remark}
\DeclareMathOperator{\Li}{Li}
\DeclareMathOperator{\FP}{FP}
\DeclareMathOperator{\CT}{CT}
\DeclareMathOperator{\Res}{Res}
\DeclareMathOperator{\Shi}{Shi}
\DeclareMathOperator{\Log}{Log}
\DeclareMathOperator{\ord}{ord}
\DeclareMathOperator{\wt}{wt}
\newcommand{\re}{\operatorname{Re}}
\newcommand{\im}{\operatorname{Im}}
\newcommand{\dd}{\,\mathrm d}
\newcommand{\Q}{\mathbb Q}
\newcommand{\C}{\mathbb C}
\newcommand{\A}{\mathcal A}
\newcommand{\NN}{\mathcal N}
\newcommand{\HH}{\mathcal H}
\newcommand{\zv}[1]{\zeta(#1)}
\pagestyle{fancy}
\fancyhf{}
\fancyhead[L]{\small Exact identities for the ProveIt project}
\fancyhead[R]{\small October 2026}
\fancyfoot[C]{\thepage}
\setlength{\headheight}{14pt}

\title{\textbf{Exact Polylogarithm Parity, Harmonic-Tail Reconstruction,\\
and Higher Polygamma Moments}\\[5pt]
\large Conjecture proofs, canonical reductions, spectral derivatives,\\
and a source-level research audit}
\author{Research continuation prepared for Vladimir Reshetnikov\\
and the ProveIt project}
\date{11 October 2026}

\begin{document}
\maketitle
\begin{abstract}
This article continues the current ProveIt polylogarithm manuscript and
the six incoming reports at the pinned source revision specified below.
We prove the optimal consecutive-moment conjecture for every pair of
adjacent even harmonic-tail orders, with an explicit nonzero rational
determinant.  An elementary implementation of established bicolored
parity proves the recorded Gaussian Family--2 identity and supplies an
all-weight family of ordinary logarithmic integrals.  For the Gamma
differential hierarchy we derive all polygamma-square coefficients and
a closed generating function for its complete digamma--trigamma sector.
A residue-cancellation theorem converts every differential polynomial
containing a polygamma derivative into an absolutely convergent
Mittag--Leffler expansion, a normalized antiderivative, and convergent
harmonic-zeta derivative coordinates.  These tools give an explicit
sixth digamma moment and constructive higher-weight formulas.
We also evaluate all integer-order quadratic Hurwitz finite parts and
their spectral jets, and thereby all products of positive argument
derivatives of generalized Stieltjes constants, keeping the original
endpoint coordinate fixed.
Classical antecedents, inherited results, new proofs, unresolved
arithmetic reductions, and reproducibility evidence are distinguished
throughout.  The remaining research questions emphasize exact
identities and proof mechanisms rather than bounds.
\end{abstract}

\vspace{3pt}
\noindent\textbf{Source baseline.}
\texttt{VladimirReshetnikov/ProveIt}.\\
Commit \texttt{f7e35be357654198c1546fd40f363694c5bcc2cc}.
The audit covers the supplied manuscript and all six incoming ZIP
inventories, with focused proof checks described in Section~\ref{audit:sec}.
The accompanying archive contains source, exact algorithms, completed
verification records, and integration notes.  No conjectured
transcendence or independence is used in a proof.

\clearpage
{\small\setlength{\parskip}{0pt}\tableofcontents}
\newpage
\section{Scope, conventions, and the proof ledger}
\label{scope:sec}

The source corpus already contains substantial all-order mathematics.
In particular, the incoming reports \cite{SourceCanonical,SourceHigher} establish Gamma-quotient harmonic
generators, arbitrary digamma-power algorithms, the fifth moment,
continuous-depth Lerch constructions, and several Stieltjes collision
formulas.  Repeating those results as new would obscure the actual
progress.  We use them as a baseline and concentrate on four concrete
tasks: resolving a recorded conjecture on an infinite family, proving a
recorded numerical polylogarithm relation, finding closed coefficient
formulas inside the higher Gamma hierarchy, and completing explicit
higher-moment and spectral-derivative calculations.

\subsection{What is proved here, and what is inherited}

\begin{center}
\small
\begin{tabularx}{\textwidth}{@{}p{.26\textwidth}X@{}}
\toprule
Topic & Precise status in this continuation \\
\midrule
Gaussian Family--2 & The numerical source equation is proved exactly,
and reduced to $\zeta(3)$, $\pi\beta(2)$, and $\pi^2\log2$.
The general parity theorem is established prior work. \\
Adjacent even tails & The first-$q$ conjecture is proved for all
$p=q-2\ge2$.  The exact determinant and reconstruction identities
are new developments relative to the inspected reports. \\
Differential normal form & The functional-monomial decomposition is
classical.  Its Gamma coefficient sectors and explicit endpoint
specialization are developed and proved here. \\
Quadratic Hurwitz products & Explicit integer finite parts and every
order-derivative jet follow from the classical product kernel, with
the local coordinate correction proved. \\
Sixth digamma moment & A complete explicit formula is supplied in
height-one Laurent slopes and one specified, convergent depth-at-most-three
zeta-derivative combination.  Its further arithmetic reduction remains open. \\
$S_6$, revised $S_8$, larger gaps & These remain conjectural or open.
Neither the parity argument nor finite numerical evidence resolves them. \\
\bottomrule
\end{tabularx}
\end{center}

The distinction between an \emph{explicit coordinate formula} and a
\emph{reduction to a prescribed constant alphabet} matters.  Defining a
constant by an unexplained divergent sum would not complete a moment
calculation.  Our remaining mixed coordinate is instead given by an
absolutely convergent logarithmic series, a finite multiple-zeta
expression, an exact value at spectral order zero, and independently
checked analytic formulas.  We do not assert that it is irreducible or
independent of the other constants.

\subsection{Harmonic, polylogarithmic, and Laurent conventions}

For $n\ge0$ and $r\ge1$,
\[
 H_n^{(r)}=\sum_{k=1}^n k^{-r},\qquad H_n=H_n^{(1)},
 \qquad H_0^{(r)}=0.
\]
For a word $\boldsymbol a=(a_1,\ldots,a_d)$ of positive integers,
\[
 H_n(\boldsymbol a)=\sum_{n\ge k_1>\cdots>k_d\ge1}
             \frac1{k_1^{a_1}\cdots k_d^{a_d}},
 \qquad H_n(\varnothing)=1.
\]
All multiple zeta functions and multiple polylogarithms in this article
use \emph{decreasing} indices:
\[
 \zeta(s,\boldsymbol a)=\sum_{n\ge1}\frac{H_{n-1}(\boldsymbol a)}{n^s},
 \qquad
 \Li_{p,q}(x,y)=\sum_{n>m\ge1}\frac{x^ny^m}{n^pm^q}.
\]
Outside their initial domains, multiple zeta functions mean their
meromorphic continuations.  Where a displayed finite combination is
holomorphic although individual terms have poles, the combination is
formed \emph{before} its Taylor coefficient is taken.

We use the conventional generalized Stieltjes constants
\begin{equation}\label{scope:stieltjes}
 \zeta(1+u,a)=\frac1u+
       \sum_{n\ge0}\frac{(-1)^n\gamma_n(a)}{n!}u^n,
 \qquad \gamma_0(a)=-\psi(a),\quad \gamma_n=\gamma_n(1).
\end{equation}
Thus the coefficient of $u$ at $a=1$ is $-\gamma_1$, not
$\gamma_1$.  A second family used below is
\begin{equation}\label{scope:lambda}
 \lambda_d=[u^1]\zeta(1+u,\{1\}^{d-1}),\qquad d\ge1.
\end{equation}
Here $[u^1]$ means the Laurent coefficient; the function may have a pole
of order $d$.  In particular $\lambda_1=-\gamma_1$.
The higher $\lambda_d$ are strict-depth Laurent slopes, not a relabeling
of the ordinary $\gamma_d$.  The local notation $\lambda_2$ for a
dilogarithm imaginary part in the source Family--2 comparison is
explicitly confined to that comparison.

The basic differential identities are
\begin{equation}\label{scope:hurwitz-differentiation}
 \partial_a\zeta(s,a)=-s\zeta(s+1,a),\qquad
 \psi^{(r)}(a)=(-1)^{r+1}r!\zeta(r+1,a)\quad(r\ge1).
\end{equation}
They follow directly from the defining series in its absolute
convergence domain and then by continuation; see also DLMF~25.11
\cite{DLMF}.  They fix the signs when translating the harmonic-tail
identities into polygamma products.

For later comparison with antiderivatives, coefficient extraction in
\eqref{scope:hurwitz-differentiation} gives the classical formulas
\begin{align}
 \gamma_n'(a)
  &=(-1)^{n+1}
     \left\{\left.\partial_s^n\zeta(s,a)\right|_{s=2}
      +n\left.\partial_s^{n-1}\zeta(s,a)\right|_{s=2}\right\},
       \label{scope:stieltjes-derivative}\\
 \int_{a_0}^{a}\gamma_n(v)\dd v
  &=\frac{(-1)^{n+1}}{n+1}
    \left.\partial_s^{n+1}
       \{\zeta(s,a)-\zeta(s,a_0)\}\right|_{s=0}.
       \label{scope:stieltjes-primitive}
\end{align}
The second term in braces in \eqref{scope:stieltjes-derivative} is
absent for $n=0$.  The derivative formula comes from
$\partial_a\zeta(1+u,a)=-(1+u)\zeta(2+u,a)$.
For the primitive use
$\partial_a\zeta(u,a)=-u\zeta(1+u,a)$, then differentiate
$n+1$ times at zero.  These formulas are included to connect the new
spectral-jet identities with the conventional Stieltjes normalization;
no priority claim is made for them.

\subsection{The endpoint coordinate is part of the definition}

If a function on $(0,1]$ has a finite principal Laurent part at zero,
we define
\begin{equation}\label{scope:finitepart}
 \FP\int_0^1 f(x)\dd x
   =\CT_{\epsilon=0}\int_\epsilon^1 f(x)\dd x.
\end{equation}
The same notation applies to finite sums of powers times logarithms,
followed by an integrable remainder.  The constant coefficient is taken
after all negative powers and all divergent logarithms are removed.
The cutoff is always the displayed coordinate $x=\epsilon$.
In particular,
\[
 \FP\int_0^1x^{-j}\dd x=-\frac1{j-1}\ (j\ge2),\qquad
 \FP\int_0^1x^{-1}(\log x)^a\dd x=0\ (a\ge0).
\]
For $h$ meromorphic near $[0,1]$ with no poles on $(0,1]$,
\begin{equation}\label{scope:boundary}
 \FP\int_0^1h'(x)\dd x=h(1)-[x^0]h(x).
\end{equation}
This is an exact coefficient identity obtained by integrating from
$\epsilon$ to $1$.  It does not replace $[x^0]h$ by the generally
nonexistent value $h(0)$.

Finally, throughout the tail-rank and differential-polynomial sections,
``rank'' and ``dimension'' refer to coefficient spaces over $\Q$.
For example, using distinct formal symbols for $\zeta(0),\zeta(2),\ldots$
does not assert their algebraic independence.  Likewise the image of a
derivation is a vector subspace, generally not an ideal.

% Self-contained section fragment for integration; ordinary amsmath/amsthm suffice.
% Define \Li as \operatorname{Li} if the host does not already do so.
\section[A bilateral proof of bicolored parity and the Gaussian Family--2 identity]{A bilateral proof of bicolored parity\\and the Gaussian Family--2 identity}
\label{bip:sec:main}

The canonical manuscript's equation \texttt{double:eq:fam2-example}
was supported by integer-relation discovery and numerical evaluation.
We give a complete proof, and simplify its right-hand side to three
single-value products.  The proof is a concrete application of known
parity theory: Nakamura's unit-circle functional relation
\cite[Proposition~1.2]{bip:Nakamura2012} precedes Panzer's general
multiple-polylogarithm parity theorem \cite[Theorem~1.3 and (3.2)]{bip:Panzer2017}.
Our contribution here is the explicit formula in the manuscript's
\emph{decreasing-index} convention, a direct bilateral proof including
indices equal to one and the collision $xy=1$, and its exact application
to the recorded Gaussian relation.  It is not a claim of priority for
unit-circle parity or for reducibility of all the constants displayed below.

\subsection{Conventions, ordinary convergence, and the finite formula}
Write
\[
 Z_{p,q}(x,y)=\Li_{p,q}(x,y)
 =\sum_{n>m\ge1}\frac{x^ny^m}{n^pm^q},\qquad p,q\ge1.
\]
In this section $|x|=|y|=1$ and $x,y\ne1$; the product $r=xy$ is
allowed to equal one.  These hypotheses include the Gaussian pairs
$(i,i)$, $(-1,i)$, $(i,-1)$, and $(i,-i)$.
The ordinary nested series converges, including when $p=1$.
Indeed, for $y\ne1$, the inner partial sums
$H_{n-1}^{(q)}(y)=\sum_{m<n}y^m/m^q$ are bounded; the sequence
$H_{n-1}^{(q)}(y)/n^p$ tends to zero and its successive differences
are absolutely summable.  Summation by parts against the bounded
partial sums of $x^n$ proves convergence.

For $|u|=1$ put
\[
 \mathcal B_j(u)=\Li_j(u)+(-1)^j\Li_j(u^{-1})\quad(u\ne1).
\]
At $u=1$ define
\[
 \mathcal B_j(1)=\begin{cases}2\zeta(j),&j\text{ even},\\0,&j\text{ odd}.
 \end{cases}
\]
In particular $\mathcal B_1(1)=0$ is a symmetric bilateral value;
it is not an assignment of a finite value to $\Li_1(1)$.
For $u=e^{2\pi i t}$, $0<t<1$, the classical inversion identity gives
\[
 \mathcal B_j(u)=-\frac{(2\pi i)^j}{j!}B_j(t),
\]
where $B_j(t)$ is the Bernoulli polynomial. Thus all these factors are
explicit rational multiples of $(i\pi)^j$ at roots of unity.

For $w=p+q$ define
\begin{align}
 \mathcal R_{p,q}(x,y)={}&(-1)^q\Biggl[
 \sum_{j=1}^q(-1)^j\binom{w-j-1}{p-1}
       \mathcal B_j(r)\Li_{w-j}(x)\notag\\
 &\qquad+\sum_{j=1}^p\binom{w-j-1}{q-1}
       \mathcal B_j(r)\Li_{w-j}(y^{-1})\notag\\
 &\qquad-\binom{w-1}{p}\Li_w(y^{-1})
       -\binom{w-1}{q}\Li_w(x)
       \notag\\
 &\qquad-\mathcal B_p(x)\Li_q(y^{-1})\Biggr]
       +(-1)^w\Li_w(r^{-1}).
 \label{bip:eq:R}
\end{align}
No $\Li_1(1)$ occurs in this expression: the arguments of the possible
order-one single values are $x$ and $y^{-1}$, both different from one;
the last single value has order $w\ge2$.

\begin{theorem}[Explicit bicolored parity, in decreasing-index convention]
\label{bip:thm:parity}
For the parameters above,
\begin{equation}
 \boxed{\quad Z_{p,q}(x,y)-(-1)^wZ_{p,q}(x^{-1},y^{-1})
      =\mathcal R_{p,q}(x,y).\quad}
 \label{bip:eq:parity}
\end{equation}
Consequently
\[
 \Re Z_{p,q}(x,y)=\frac12\mathcal R_{p,q}(x,y)
       \quad(w\text{ odd}),
 \qquad
 \Im Z_{p,q}(x,y)=\frac1{2i}\mathcal R_{p,q}(x,y)
       \quad(w\text{ even}).
\]
At roots of unity these are finite formulas in single polylogarithms
and powers of $\pi$.  The opposite parity component is not evaluated
by this theorem.
\end{theorem}

\subsection{Proof with the summation order made explicit}
For $k\ge1$ form the absolutely convergent bilateral sum
\[
 C_{p,q}(k;r)=\sum_{m\in\mathbb Z\setminus\{0,-k\}}
       \frac{r^m}{(m+k)^pm^q}.
\]
The denominator is $O(|m|^w)$ as $|m|\to\infty$, and $w\ge2$.
The elementary partial-fraction identity
\begin{align}
 \frac1{(m+k)^pm^q}
 ={}&\sum_{j=1}^q(-1)^{q-j}\binom{w-j-1}{p-1}
       \frac1{k^{w-j}m^j}\notag\\
 &+(-1)^q\sum_{j=1}^p\binom{w-j-1}{q-1}
       \frac1{k^{w-j}(m+k)^j}
 \label{bip:eq:pf}
\end{align}
holds away from its poles.  It follows either by matching the principal
parts at $m=0,-k$, after which the difference is a rational function
with no poles and vanishing at infinity, or by induction on $p+q$.

Sum \eqref{bip:eq:pf} over a symmetric cutoff in $m$.  The possible
$j=1$ terms are legitimate in this order:
\[
 \underset{M\to\infty}{\lim}\sum_{0<|m|\le M}\frac{r^m}{m^j}
      =\mathcal B_j(r).
\]
For $r\ne1$, the two one-sided $j=1$ series converge by Dirichlet's
test; for $r=1$ they cancel at every symmetric cutoff.  Translating
the cutoff by the fixed integer $k$ changes the sum by finitely many
terms tending to zero, so the substitution $m+k=\ell$ is valid.
Removing the omitted terms $m=-k$ and $\ell=k$ gives
\begin{align}
 C_{p,q}(k;r)={}&
 \sum_{j=1}^q(-1)^{q-j}\binom{w-j-1}{p-1}
   \frac{\mathcal B_j(r)}{k^{w-j}}\notag\\
 &+(-1)^qr^{-k}\sum_{j=1}^p\binom{w-j-1}{q-1}
   \frac{\mathcal B_j(r)}{k^{w-j}}\notag\\
 &-\frac{(-1)^q}{k^w}
       \left[\binom{w-1}{p}r^{-k}+\binom{w-1}{q}\right].
 \label{bip:eq:C}
\end{align}
The last line uses the elementary hockey-stick sums
$\sum_{j=1}^q\binom{w-j-1}{p-1}=\binom{w-1}{p}$ and its
$p,q$ counterpart.

For $0<\rho<1$ multiply \eqref{bip:eq:C} by $\rho^kx^k$ and sum over
$k\ge1$. All the original double sums are now absolutely convergent.
On the right, let $\rho\uparrow1$ in the resulting single
polylogarithms; their arguments are $x$ and $y^{-1}$, neither equal
to one.  The result is
\begin{align}
 \mathcal A={}&(-1)^q\Biggl[
 \sum_{j=1}^q(-1)^j\binom{w-j-1}{p-1}
       \mathcal B_j(r)\Li_{w-j}(x)\notag\\
 &\quad+\sum_{j=1}^p\binom{w-j-1}{q-1}
       \mathcal B_j(r)\Li_{w-j}(y^{-1})
       -\binom{w-1}{p}\Li_w(y^{-1})
       -\binom{w-1}{q}\Li_w(x)\Biggr].
 \label{bip:eq:A}
\end{align}

To identify the left side in the ordinary nested convention, split
its lattice points, with $n=m+k$, into three disjoint regions:
$n>m>0$, $0>n>m$, and $n>0>m$. They yield, as $\rho\uparrow1$,
\begin{equation}
 \mathcal A=Z_{p,q}(x,y)+(-1)^wZ_{q,p}(y^{-1},x^{-1})
            +(-1)^q\Li_p(x)\Li_q(y^{-1}).
 \label{bip:eq:regions}
\end{equation}
For completeness, the gap Abel factors $\rho^{n-m}$ used here need
not be identified informally with outer-index radial factors.  The
first region is $Z_{p,q}(\rho x,y/\rho)$ and has the integral
\begin{equation}
 Z_{p,q}(a,b)=\frac1{(p-1)!}
    \int_0^1(-\log t)^{p-1}
            \frac{a\Li_q(abt)}{1-at}\,dt.
 \label{bip:eq:integral}
\end{equation}
For $a=\rho x$, $ab=r$, its limit follows by dominated convergence:
$1-\rho xt$ is bounded away from zero uniformly for $\rho$ near one;
$\Li_q(rt)$ is bounded when $q>1$ and at worst logarithmic near
$t=1$ when $q=1,r=1$.  These are integrable majorants.  The same
integral identifies this limit with the ordinary nested series via
outer-index Abel summation.  The negative region is handled
identically, and the cross region is the product
$(-1)^q\Li_p(\rho x)\Li_q(\rho/y)$.  This proves
\eqref{bip:eq:regions} for all stated boundary cases.

Finally the convergent stuffle identity gives
\[
 Z_{q,p}(y^{-1},x^{-1})=
 \Li_q(y^{-1})\Li_p(x^{-1})-\Li_w(r^{-1})
 -Z_{p,q}(x^{-1},y^{-1}).
\]
It follows by multiplying finite partial sums and passing to their
limits; both nested series have already been shown to converge.
Substituting this into \eqref{bip:eq:regions} and using
$\mathcal B_p(x)=\Li_p(x)+(-1)^p\Li_p(x^{-1})$ gives
\eqref{bip:eq:parity}.  Complex conjugation proves the two displayed
real/imaginary consequences.

The collision $r=1$ has therefore been proved directly. One can also
see why it is harmless in \eqref{bip:eq:R}: its $\mathcal B_1(r)$
coefficient is a constant multiple of
$\Li_{w-1}(y^{-1})-\Li_{w-1}(x)$, which vanishes at $r=1$.
The jump in the individual bilateral order-one factor cannot create
a jump in the complete expression.

\subsection{Proof and simplification of the recorded Family--2 identity}
Let $L=\log2$, $G=\beta(2)$, and $v=(1+i)/2$. The exact depth-one
values required at weight three are
\[
 \Li_1(i)=-\frac L2+\frac{i\pi}4,\quad
 \Li_2(i)=-\frac{\pi^2}{48}+iG,\quad
 \Li_3(i)=-\frac{3\zeta(3)}{32}+\frac{i\pi^3}{32},
\]
with $\Li_n(-i)=\overline{\Li_n(i)}$ and
$\Li_n(-1)=-(1-2^{1-n})\zeta(n)$ for $n>1$.
Here $\mathcal B_1(-1)=0$ and
$\mathcal B_2(-1)=-\pi^2/6$. Inserting these numbers in
\eqref{bip:eq:R}, first at $(p,q)=(1,2)$ and then at $(2,1)$,
gives
\begin{align}
 \Re\Li_{1,2}(i,i)
   &=\frac{33}{64}\zeta(3)-\frac\pi4G+\frac{\pi^2L}{24},
       \label{bip:eq:12}\\
 \Re\Li_{2,1}(i,i)
   &=\frac{15}{64}\zeta(3)-\frac{\pi^2L}{32}.
       \label{bip:eq:21}
\end{align}
Consequently the desired difference has the short exact form
\begin{equation}
 \boxed{\quad
 \Re\bigl(\Li_{1,2}(i,i)-\Li_{2,1}(i,i)\bigr)
 =\frac9{32}\zeta(3)-\frac\pi4G+\frac7{96}\pi^2\log2.
 \quad}
 \label{bip:eq:family2-short}
\end{equation}

To identify this with the exact equation printed in the source, write
$\mu_2=\Re\Li_2(v)$, $\lambda_2=\Im\Li_2(v)$ and
$h_{a,b}=\Re\Li_{a,b}(i,1)$. The classical Landen formula
\[
 \Li_2\!\left(\frac z{z-1}\right)
      =-\Li_2(z)-\frac12\log^2(1-z)
\]
(applied along the segment $z=it$, $0\le t\le1$) gives
\[
 \mu_2=\frac{5\pi^2}{96}-\frac{L^2}{8},\qquad
 \lambda_2=G-\frac{\pi L}{8}.
\]
Its derivative is an identity of elementary logarithms and its
constant is fixed at $z=0$, so no numerical recognition is needed.
The shuffle identity immediately gives
\[
 h_{1,1}=\frac{L^2}{8}-\frac{\pi^2}{32}.
\]
For the remaining height-one value put
\[
 S(z)=\Li_{2,1}(z,1)=\frac12\int_0^z\frac{\log^2(1-t)}t\,dt.
\]
The needed inversion is already recorded after the proof of the source's
Proposition~1.12, ``Nielsen inversion with explicit endpoint data''
\cite{SourceManuscript}; both the inversion and its Gaussian specialization
are also explicit in Panzer's discussion following Corollary~1.4
\cite{bip:Panzer2017}.  To make the present proof independent of an
unlabelled source display, let $u=\log(-z)$ on
$\mathbb C\setminus[0,\infty)$.  The identity
$\log(1-1/z)=\log(1-z)-u$ shows by differentiation that
\[
 S(z)+S(1/z)-\Li_3(z)+u\Li_2(z)+\frac{u^3}{6}
\]
is constant.  Here
\[
 S(1)=\frac12\int_0^1\frac{\log^2t}{1-t}\,dt=\zeta(3),
\]
where the last equality follows by integrating the geometric series.
Letting $z\to1$ along the upper unit semicircle gives $u\to-i\pi$;
the two cubic imaginary terms cancel because $\zeta(2)=\pi^2/6$.
The constant is therefore $\zeta(3)$, proving
\begin{equation}\label{bip:eq:nielsen-inversion}
 S(z)+S(1/z)=\Li_3(z)-\log(-z)\Li_2(z)
               -\frac16\log^3(-z)+\zeta(3).
\end{equation}
At $z=i$, conjugation gives $S(1/i)=\overline{S(i)}$ and
$\log(-i)=-i\pi/2$, so
\[
 h_{2,1}=\frac{29}{64}\zeta(3)-\frac\pi4G.
\]
Finally $\Li_1(i)\Li_2(i)=\Li_{1,2}(i,1)+2\Li_{2,1}(i,1)$
gives
\[
 h_{1,2}=\Re\bigl(\Li_1(i)\Li_2(i)\bigr)-2h_{2,1}
       =\frac{\pi^2L}{96}+\frac\pi4G-\frac{29}{32}\zeta(3).
\]
Direct rational substitution now yields
\[
 3\zeta(3)-4L\mu_2-\pi\lambda_2-4Lh_{1,1}+3h_{1,2}
 =\frac9{32}\zeta(3)-\frac\pi4G+\frac7{96}\pi^2L.
\]
This proves \texttt{double:eq:fam2-example} precisely as written,
and proves its stronger simplification \eqref{bip:eq:family2-short}.
The statement needs no conjectured independence of the three constants.

\subsection{Mixed-color examples beyond the displayed source equation}
The same finite formula gives the following exact evaluations:
\begin{align}
 \Re\bigl(\Li_{1,2}(-1,i)-\Li_{2,1}(i,-1)\bigr)
   &=\frac{15}{16}\zeta(3)-\frac\pi2G+\frac{\pi^2L}{48},
       \label{bip:eq:mix3}\\
 \Im\bigl(\Li_{1,3}(-1,i)-\Li_{3,1}(i,-1)\bigr)
   &=\frac{3\pi^3L}{32}-3\beta(4)+\frac{21\pi\zeta(3)}{64},
       \label{bip:eq:mix4}\\
 \Re\bigl(\Li_{1,4}(-1,i)-\Li_{4,1}(i,-1)\bigr)
   &=\frac{7\pi^4L}{11520}-\frac\pi2\beta(4)
          +\frac{\pi^2\zeta(3)}{32}+\frac{135}{128}\zeta(5),
       \label{bip:eq:mix5}\\
 \Im\bigl(\Li_{1,5}(-1,i)-\Li_{5,1}(i,-1)\bigr)
   &=\frac{5\pi^5L}{512}-5\beta(6)
          +\frac{465\pi\zeta(5)}{1024}
          +\frac{3\pi^3\zeta(3)}{64}.
       \label{bip:eq:mix6}
\end{align}
Both indices \emph{and} colors are swapped in these antisymmetric
expressions. They are elementary rational specializations of
Theorem~\ref{bip:thm:parity}, not additional parity conjectures.
The accompanying symbolic generator derives the Gaussian formulas at
all requested weights; a separate Mellin-integral evaluator checks the
left sides without reusing the theorem as its evaluation algorithm.

As a collision check, the case $p=q=1$, $(x,y)=(i,-i)$ gives
\[
 \Im\Li_{1,1}(i,-i)=G-\frac{\pi\log2}{8}.
\]
Here the symmetric bilateral value $\mathcal B_1(1)=0$ is essential to
avoid introducing a spurious divergent $\Li_1(1)$.


\subsection{An all-weight family of ordinary logarithmic integrals}
The single-inner-index specialization has an additional cancellation:
no even beta values remain in its real odd-weight component.

\begin{corollary}[Odd logarithmic moments]
For $m\ge1$, put $w=2m+1$ and
\[
 I_m=\frac1{(2m-1)!}\int_0^1
       \frac{x(-\log x)^{2m-1}\log(1+x)}{1+x^2}\,dx.
\]
Then
\begin{align}
 I_m={}&\frac12(1-2^{-2m})
        \left(1-\frac{2m+1}{2^{2m+1}}\right)\zeta(2m+1)\notag\\
 &-\sum_{k=1}^{m-1}
  \frac{(1-2^{1-2k})(1-2^{1-w+2k})}{2^{w-2k}}
       \zeta(2k)\zeta(w-2k)\notag\\
 &-\frac12(1-2^{1-2m})(1-2^{-2m})\zeta(2m)\log2.
 \label{bip:eq:moment-family}
\end{align}
In particular,
\begin{align*}
 I_1&=\frac{15}{64}\zeta(3)-\frac{\pi^2\log2}{32},\\
 I_2&=\frac{405}{1024}\zeta(5)-\frac{\pi^2\zeta(3)}{128}
                     -\frac{7\pi^4\log2}{1536},\\
 I_3&=\frac{7623}{16384}\zeta(7)-\frac{5\pi^2\zeta(5)}{2048}
       -\frac{7\pi^4\zeta(3)}{7680}-\frac{31\pi^6\log2}{61440}.
\end{align*}
\end{corollary}
\begin{proof}
Equation~\eqref{bip:eq:integral} with $(x,y)=(i,i)$ and $q=1$
gives $I_m=\Re Z_{2m,1}(i,i)$, since
$\Li_1(-t)=-\log(1+t)$ and $\Re(i/(1-it))=-t/(1+t^2)$.
In \eqref{bip:eq:R} use
$\mathcal B_{2k}(-1)=2\Li_{2k}(-1)$ and
$\mathcal B_{2k+1}(-1)=0$. Writing
$u_j=\Re\Li_j(i)$, the real part simplifies before any
zeta evaluation to
\[
 2I_m=-2\sum_{k=1}^{m}\Li_{2k}(-1)u_{w-2k}
           +w u_w-u_{2m}\log2-\Li_w(-1).
\]
Now $u_1=-\log2/2$ and
$u_j=-2^{-j}(1-2^{1-j})\zeta(j)$ for $j\ge2$.
Separate the $k=m$ term and substitute these expressions to obtain
\eqref{bip:eq:moment-family}.  Euler's even-zeta evaluations yield
the three displayed cases.
\end{proof}

Two weight-three integral identities also follow directly:
\begin{align*}
 \int_0^1\frac{x\log x\log(1+x)}{1+x^2}\,dx
  &=\frac{\pi^2\log2}{32}-\frac{15}{64}\zeta(3),\\
 \int_0^1\frac{x\Li_2(-x)}{1+x^2}\,dx
  &=\frac\pi4G-\frac{\pi^2\log2}{24}-\frac{33}{64}\zeta(3).
\end{align*}
These are ordinary real integrals with their constants fixed by the
convergent endpoint representation. They may be useful as simple
benchmark identities alongside the bicolored-polylogarithm form.

\subsection{Audit conclusion and what this does not settle}
The numeric Family--2 example can now be marked proved and replaced,
if desired, by \eqref{bip:eq:family2-short}; its older coordinates can
be preserved as an exact equivalent identity. No false equality was
found in that example. The correction is to its evidence status and
to an unnecessarily large choice of constants.

The master identity proves the real components in odd weight and the
imaginary components in even weight for every admissible Gaussian
mixed pair, and provides the same statement at all roots of unity.
It therefore supplies an exact replacement for searches in these
parity components. It does not prove nonreduction in the complementary
components.
In particular $S_6$ and $S_8$ have odd total weights seven and nine
and are imaginary components. They lie in the opposite parity sector;
this theorem proves neither frozen candidate. The rational
weight-nine vector already excluded in the source remains false, and
the newer $S_6,S_8$ candidates remain conjectural. The higher golden
ladders are also unaffected.

\paragraph{Attribution safeguard.}
When comparing formulas, Panzer uses increasing summation indices,
so both indices and arguments must be reversed to use his notation.
Panzer's footnote~16 explicitly flags a transposed product in
Nakamura's printed equation~(1.3): the product
$\zeta(a;-y)\zeta(b;x)$ must be replaced by
$\zeta(a;x)\zeta(b;-y)$ in that convention. The direct proof above
rebuilds the products from the three lattice regions and therefore
avoids relying on the misprinted equation. This is a literature
normalization issue, not a newly discovered source-manuscript error.

% Bibliography entries, if not already present in the host:
% \bibitem{bip:Panzer2017}
% E. Panzer, The parity theorem for multiple polylogarithms,
% Journal of Number Theory 172 (2017), 93--113.
% \url{https://arxiv.org/abs/1512.04482}.
% \bibitem{bip:Nakamura2012}
% T. Nakamura, A simple proof of the functional relation for the Lerch type
% Tornheim double zeta function, Tokyo Journal of Mathematics 35 (2012),
% 333--337. \url{https://arxiv.org/abs/1012.1144}.

% Research fragment for the parent article.  Not a standalone document.
% Authored against ProveIt f7e35be357654198c1546fd40f363694c5bcc2cc.
% Source conjecture: incoming Canonical Digamma and Fifth Harmonic Bridge,
% sections/09_questions.tex, Conjecture conj:firstq.
% Source exact tail formulas: sections/07_tail_rank.tex, tr:Aentries,
% tr:forward, tr:backward, tr:operator.

\section{The optimal moment prefix for adjacent even tail orders}
\label{adj:section}

The incoming \emph{Canonical Digamma Polynomials and the Fifth Harmonic
Bridge} proves that finitely many initial centered moments generate the
full formal coordinate space of two even harmonic tails.  For equal
orders it proves the optimal consecutive prefix.  Its
Conjecture~\texttt{conj:firstq} asks for that optimal prefix for every
unequal pair.  We prove the entire first off-diagonal family:
\[
 p=2b-2,\qquad q=2b,\qquad b\ge2.
\]
The proof also gives an exact determinant with a positive rational
recurrence, and hence a family of ordinary convergent identities for
products of higher polygammas.  The result concerns rational coefficient
rank.  It does not assert arithmetic independence of the constants that
those coefficients multiply.

\subsection{The inherited exact coordinate system}

Write
\[
 A_p(n)=\zeta(p)-H_n^{(p)}=\zeta(p,n+1),\qquad
 x_n=n+\tfrac12,
\]
and let $D_r$ be the continued value at $s=-2r$ of
\[
 D_{p,q}(s)=\sum_{n=0}^{\infty}A_p(n)A_q(n)x_n^{-s}.
\]
For even $p,q$, these particular continued values are regular.  The
source tail theorem supplies numbers $C_j$ satisfying the inverse
Bernoulli transforms
\begin{align}
 D_r&=\sum_{j=0}^{r}
 \frac{\binom{2r}{2j}}{2j+1}B_{2r-2j}(1/2)C_j,
 \label{adj:forward}\\
 C_j&=\sum_{r=0}^{j}\binom{2j+1}{2r}2^{-2(j-r)}D_r.
 \label{adj:backward}
\end{align}
For $q=2b$ and $j\ge b$, they have the finite coordinate expansion
\begin{equation}\label{adj:block}
 C_j=\sum_{k=0}^{b-1}a_{jk}\zeta(2k),\qquad \zeta(0)=-\tfrac12,
\end{equation}
where, putting $d=2j+1$ and $e=d-p-q+2k+1$,
\begin{equation}\label{adj:a}
 a_{jk}=B_e\left\{
 \frac{\binom{d-q+1}{p-2k}}{d-q+1}
 +\frac{\binom{d-p+1}{q-2k}}{d-p+1}
 \right\}\quad(e\ge0),
\end{equation}
and $a_{jk}=0$ if $e<0$.  Each binomial summand in
\eqref{adj:a} is omitted unless its lower index is at least $2$ and
no greater than its upper index.  In particular, a would-be lower
index zero is \emph{not} retained.  These formulas are prior results
of the inspected reports; the new argument is the nonvanishing theorem
below.

For reference, the finite mechanism behind the inherited formulas is
summation by parts with
\[
 \sum_{n=0}^{N-1}x_n^{2r}
 =\sum_{j=0}^{r}\frac{\binom{2r}{2j}}{2j+1}
 B_{2r-2j}(1/2)N^{2j+1},
 \qquad A_p(k-1)=A_p(k)+k^{-p}.
\]
The resulting ordered finite sums reduce by Faulhaber polynomials.
Their constant coefficients give \eqref{adj:block}--\eqref{adj:a};
all subtracted nonnegative even powers have continued sum zero because
$\zeta(-2m,1/2)=0$.  Thus the conventions in
\eqref{adj:forward} are tied to the original centered cutoff, rather
than to an unspecified regularization.

\subsection{A universal positive recurrence}

Let
\[
 h(t)=\coth(t/2)=\sum_{k=0}^{\infty}h_k t^{2k-1},\qquad
 h_k=\frac{2B_{2k}}{(2k)!},\qquad h_0=2.
\]
The ordinary Bernoulli sign formula gives
\begin{equation}\label{adj:hsign}
 h_k=(-1)^{k-1}\beta_k,\qquad
 \beta_k=\frac{2|B_{2k}|}{(2k)!}
 =\frac{4\zeta(2k)}{(2\pi)^{2k}}>0\quad(k\ge1).
\end{equation}
This is the $x=0$ case of the classical Bernoulli Fourier formula
(DLMF~24.8.1).  Define rational numbers $c_n$ by
\begin{align}
 c_1&=1,\label{adj:c1}\\
 c_n&=\frac{1}{4(2n-1)^2}
 \sum_{r=1}^{n-1}\beta_{n-r}
 \left\{2r(2r-1)+(2n-1)(2n-2)-2\mathbf1_{r=1}\right\}c_r
 \quad(n\ge2).
 \label{adj:positive-recurrence}
\end{align}
Every coefficient in this recurrence is strictly positive.  Indeed,
for $r=1$ the expression in braces is $(2n-1)(2n-2)>0$, and for
$r\ge2$ both its displayed summands are positive.  Therefore
\begin{equation}\label{adj:cpositive}
 c_n\in\mathbb Q_{>0}\quad(n\ge1).
\end{equation}
The first values are
\[
 c_1=1,\quad c_2=\frac1{36},\quad c_3=\frac{11}{5400},\quad
 c_4=\frac{113}{705600},\quad
 c_5=\frac{2509}{190512000},\quad
 c_6=\frac{309713}{276623424000}.
\]

\begin{theorem}[Optimal prefix and exact determinant for adjacent orders]
\label{adj:main}
Let $b\ge2$, $p=2b-2$, and $q=2b$.  Define the square late-block
matrix
\[
 M_b=(a_{b+r,k})_{0\le r,k<b}
\]
by \eqref{adj:a}.  Then
\begin{equation}\label{adj:determinant}
 \det M_b=
 (-1)^{b(b-1)/2}\frac{((2b+1)!!)^2}{b!}\,c_{b+1}\ne0.
\end{equation}
Consequently $D_0,\ldots,D_{q-1}$ have full formal rank $q$ in
the coordinate system
\[
 C_0,\ldots,C_{b-1},\quad\zeta(0),\zeta(2),\ldots,\zeta(2b-2).
\]
Every later $D_r$ is an explicit rational combination of these first
$q$ moments.  In this coordinate ordering, the determinant of the
complete first-$q$ moment matrix is
\begin{equation}\label{adj:full-determinant}
 (-1)^{b(b-1)/2}
 \frac{((2b+1)!!)^2c_{b+1}}{b!(4b-1)!!}.
\end{equation}
\end{theorem}

\begin{proof}
For an even polynomial
\[
 Q(t)=\sum_{n=1}^{b}q_n t^{2n},
\]
introduce the operator
\begin{equation}\label{adj:operator}
 \mathcal L_Q(t)=\bigl(Q''(t)-Q''(0)\bigr)h(t)+(Q(t)h(t))''.
\end{equation}
It is an odd formal power series, and is holomorphic at zero when
regarded analytically near zero.  The subtraction of $Q''(0)$ is
essential: it is precisely the omitted lower-index-zero term in
\eqref{adj:a}.

Let $K_b$ be the matrix mapping $(q_1,\ldots,q_b)$ to the
coefficients of $t,t^3,\ldots,t^{2b-1}$ in $\mathcal L_Q$.
Direct coefficient extraction from \eqref{adj:a} gives
\begin{equation}\label{adj:matrix-conjugacy}
 K_b=\operatorname{diag}_{0\le r<b}\left(\frac2{(2r+1)!}\right)
 M_b^{\mathrm{rev}}
 \operatorname{diag}_{1\le n\le b}\bigl((2n)!\bigr),
\end{equation}
where $M_b^{\mathrm{rev}}$ reverses the column order.  This can
also be read off from the inherited identity
$2f^{(2b)}=\mathcal L_Q$ for the column generators, but no
analytic continuation is needed for the coefficient calculation.
More explicitly, with row $r=N-2$ and column $n$,
\[
 (K_b)_{N-2,n-1}
 =h_{N-n}\{2n(2n-1)+(2N-1)(2N-2)-2\mathbf1_{n=1}\}
\]
when $n\le N$, and the entry is zero when $n>N$.

Write $Q_*(t)=\sum_{n\ge1}q_n^*t^{2n}$ for the formal solution
of $\mathcal L_{Q_*}=0$ normalized by $q_1^*=1$.  For $n\ge2$,
the coefficient of $t^{2n-3}$ gives
\begin{equation}\label{adj:q-recurrence}
 4(2n-1)^2q_n^*
 +\sum_{r=1}^{n-1}h_{n-r}
 \left\{2r(2r-1)+(2n-1)(2n-2)-2\mathbf1_{r=1}\right\}q_r^*=0.
\end{equation}
The nonzero coefficient of $q_n^*$ proves existence and uniqueness
recursively.  Equations \eqref{adj:hsign} and
\eqref{adj:positive-recurrence} show that
\begin{equation}\label{adj:qsign}
 q_n^*=(-1)^{n-1}c_n.
\end{equation}
In particular no coefficient vanishes.

The matrix $K_b$ is upper Hessenberg.  Its row corresponding to
$n=2,\ldots,b+1$ has the nonzero next-column entry
$d_n=4(2n-1)^2$ whenever that column is present.  Successively
using these entries to eliminate $q_2,\ldots,q_b$ shows
\begin{equation}\label{adj:hessenberg-det}
 \det K_b=(-1)^b q_{b+1}^*\prod_{n=2}^{b+1}d_n
 =4^b((2b+1)!!)^2c_{b+1}.
\end{equation}
Here the last constraint is the one that would determine
$q_{b+1}^*$ in the infinite formal solution; a polynomial has
$q_{b+1}=0$.  One explicit determinant argument replaces the first
column by itself plus $q_2^*,\ldots,q_b^*$ times the successive
columns.  Its first $b-1$ entries become zero and its last entry
is $-d_{b+1}q_{b+1}^*$.  Expanding along that column leaves a
lower triangular minor with diagonal $d_2,\ldots,d_b$, proving
\eqref{adj:hessenberg-det} with its stated sign.
This also provides a determinant-free proof of
injectivity: the first $b-1$ constraints force
$q_n=q_1q_n^*$ through $n=b$, and the final one forces
$q_1q_{b+1}^*=0$, hence $Q=0$.

Taking determinants in \eqref{adj:matrix-conjugacy}, and using
\[
 \prod_{n=1}^{b}\frac{(2n-1)!}{(2n)!}=\frac1{2^bb!},
\]
proves \eqref{adj:determinant}.  The transform
\eqref{adj:forward} is lower triangular with diagonal
$1/(2r+1)$, and the first $b$ block rows are the identity on
$C_0,\ldots,C_{b-1}$.  Its product of diagonal entries through
$r=2b-1$ is $1/(4b-1)!!$.  This proves
\eqref{adj:full-determinant} and the rank assertion.  Finally
$M_b^{-1}$ reconstructs every even-zeta coordinate from
$C_b,\ldots,C_{2b-1}$, while \eqref{adj:backward} obtains all
these blocks from the first $2b$ moments.  Equations
\eqref{adj:a} and \eqref{adj:forward} then reconstruct every
later moment by finite rational arithmetic.
\end{proof}

\paragraph{A clarification about the range of the theorem.}
The determinant formula proves the full infinite family $q-p=2$.
The case $p=q$ was already proved in the incoming report.  General
unequal gaps $q-p\ge4$ remain open here; the argument is not a
proof of the complete source conjecture.  Its positivity mechanism
uses a single free initial coefficient, which is specific to the
second-order operator in \eqref{adj:operator}.
Isolated larger-gap pairs, including $(p,q)=(2,6)$, were already
treated in the source; the remaining question is the uniform assertion.

\subsection{Ordinary convergent harmonic and polygamma identities}

The inverse matrix produces explicit series identities without an
unevaluated finite-part symbol.  For $0\le k<b$, define
\begin{align}
 w_{b,k,r}
 &=\sum_{t=0}^{b-1}(M_b^{-1})_{kt}
 \binom{2(b+t)+1}{2r}2^{-2(b+t-r)},
 \label{adj:w}\\
 P_{b,k}(x)&=\sum_{r=0}^{2b-1}w_{b,k,r}x^{2r},
 \label{adj:P}
\end{align}
where a summand in \eqref{adj:w} is zero if $r>b+t$.  Put
\begin{equation}\label{adj:alpha}
 a_0=\frac1{(2b-3)(2b-1)},\qquad
 a_1=-\frac{2b-2}{24(2b-1)}-\frac{2b}{24(2b-3)}.
\end{equation}

\begin{corollary}[All even-zeta coordinates from adjacent tails]
\label{adj:ordinary}
For every $b\ge2$ and $0\le k<b$,
\begin{equation}\label{adj:ordinary-family}
 \begin{split}
 \sum_{n=0}^{\infty}\bigl\{&P_{b,k}(x_n)A_{2b-2}(n)A_{2b}(n)
       -a_0w_{b,k,2b-1}x_n^2\\
 &\hspace{12mm}-a_0w_{b,k,2b-2}-a_1w_{b,k,2b-1}\bigr\}
 =\zeta(2k).
 \end{split}
\end{equation}
The series is absolutely convergent.  Replacing the tail product by
\begin{equation}\label{adj:polygamma}
 A_{2b-2}(n)A_{2b}(n)
 =\frac{\psi^{(2b-3)}(n+1)\psi^{(2b-1)}(n+1)}
 {(2b-3)!(2b-1)!}
\end{equation}
gives the corresponding ordinary polygamma-product identity.
\end{corollary}

\begin{proof}
By \eqref{adj:block}, inversion, and \eqref{adj:backward},
$\sum_rw_{b,k,r}D_r=\zeta(2k)$.  The centered Hurwitz
expansion gives
\[
 A_{2b-2}(n)A_{2b}(n)
 =a_0x_n^{4-4b}+a_1x_n^{2-4b}+O(x_n^{-4b}).
\]
The largest possible degree in \eqref{adj:P} is $4b-2$.
Thus the two displayed subtractions remove every nondecaying term,
and the summand is $O(x_n^{-2})$.  Centered subtraction computes
the same continued moment because the removed monomials have
$\zeta(0,1/2)=\zeta(-2,1/2)=0$.  This proves the identity and
absolute convergence.  Formula \eqref{adj:polygamma} is the
standard Hurwitz-zeta expression for odd-order polygammas.
\end{proof}

\subsection{Two exact rational evaluations}

For $(p,q)=(4,6)$, the late block matrix and its inverse are
\begin{equation}\label{adj:46matrix}
 M_3=\begin{pmatrix}
 0&3/4&1/4\\5/12&2/3&-1/12\\1&-3/8&1/12
 \end{pmatrix},\qquad
 M_3^{-1}=\frac1{339}
 \begin{pmatrix}-28&180&264\\136&288&-120\\948&-864&360\end{pmatrix}.
\end{equation}
In particular $\det M_3=-113/384$.  Define
\begin{equation}\label{adj:P46}
 \begin{split}
 P_{4,6}(x)={}&67+15516x^2+324320x^4+1434496x^6\\
             &+1601280x^8+371712x^{10}.
 \end{split}
\end{equation}
Then the first row of \eqref{adj:46matrix} after the inverse
Bernoulli transform gives the ordinary identity
\begin{equation}\label{adj:46rational}
 \sum_{n=0}^{\infty}\left\{
 P_{4,6}(x_n)A_4(n)A_6(n)
 -\frac{123904}{5}x_n^2-\frac{316928}{5}
 \right\}=-21696.
\end{equation}
The right side is $43392\zeta(0)$.  Thus every double-zeta
coordinate and every positive even-zeta coordinate cancels
exactly.

For $(p,q)=(6,8)$, $\det M_4=2509/5120$.  Define
\begin{equation}\label{adj:P68}
 \begin{split}
 P_{6,8}(x)={}&3705+880180x^2+32541904x^4+389667904x^6\\
 &+1652632320x^8+2454887424x^{10}\\
 &+1189933056x^{12}+138854400x^{14}.
 \end{split}
\end{equation}
Its rational-coordinate identity is
\begin{equation}\label{adj:68rational}
 \sum_{n=0}^{\infty}\left\{
 P_{6,8}(x_n)A_6(n)A_8(n)
 -\frac{27770880}{7}x_n^2-\frac{98910208}{5}
 \right\}=-7707648.
\end{equation}
The product in \eqref{adj:46rational} is
$\psi^{(3)}(n+1)\psi^{(5)}(n+1)/720$; the product in
\eqref{adj:68rational} is
$\psi^{(5)}(n+1)\psi^{(7)}(n+1)/604800$.
The coefficient files also contain the corresponding exact
$\zeta(2),\zeta(4)$, and $\zeta(6)$ rows where applicable.

\subsection{A hypergeometric expression for the determinant generator}

The formal solution used in the proof has a separate exact expression.
It is not needed for nonvanishing, but supplies a second coefficient
algorithm.  Dividing \eqref{adj:operator} by $2h(t)$ gives
\begin{equation}\label{adj:Qode}
 Q_*''(t)-\operatorname{csch}(t)Q_*'(t)
 +\frac{Q_*(t)}{4\sinh^2(t/2)}=1.
\end{equation}
Set $z=\tanh(t/2)$ and $Q_*(t)=4z^2F(z^2)$.  If
$F(u)=\sum_{n\ge0}f_nu^n$, then
\begin{equation}\label{adj:frecurrence}
 f_0=1,\qquad (n+\tfrac12)^2f_n-n^2f_{n-1}=\tfrac14
 \quad(n\ge1),
\end{equation}
and hence
\begin{equation}\label{adj:fclosed}
 f_n=\left(\frac{n!}{(3/2)_n}\right)^2
 \sum_{k=0}^{n}\left(\frac{(1/2)_k}{k!}\right)^2.
\end{equation}
This follows by dividing \eqref{adj:frecurrence} by its homogeneous
product and telescoping.  It is also a direct formal derivation of
\eqref{adj:Qode}; the series converges near the origin, since the
coefficients in \eqref{adj:fclosed} have at most polynomial growth.
Consequently
\begin{equation}\label{adj:detgenerator}
 \sum_{n=1}^{\infty}(-1)^{n-1}c_nt^{2n}
 =4\tanh^2(t/2)
 \sum_{n=0}^{\infty}f_n\tanh^{2n}(t/2)
\end{equation}
as an analytic identity near zero.  The finite binomial-square sum
in \eqref{adj:fclosed} connects the universal determinant
coefficients to the classical complete-elliptic-integral
hypergeometric series.  The present theorem uses this connection as
an alternative exact representation, without claiming a new
historical priority for that classical series.

\subsection{Checks, source audit, and the remaining conjecture}

The accompanying \texttt{verify\_adjacent\_tail\_prefix.py} compares
the original Bernoulli-binomial matrix with an independent
differential-operator construction, checks the determinant formula
through $b=16$, and independently reconstructs the two displayed
examples from the source $U/R$ block formulas, retaining formal
double-zeta coordinates until their exact cancellation.  Its
numerical option evaluates rational and nonrational rows by direct
tail sums followed by an independently generated centered-Hurwitz
tail.  Those residuals are diagnostics, not interval certificates
or substitutes for the proofs above.

No mathematical error in the inherited tail theorem is needed for
this advance.  Two convention errors are explicitly excluded:
omitting $Q''(0)h(t)$ changes the operator, and omitting the column
reversal changes the determinant sign.  The source optimal-prefix
conjecture should be updated to record the newly proved adjacent
family, while retaining open status for general gaps at least four.

\section[Gamma differential polynomials and exact finite-part reduction]{Gamma differential polynomials\\and exact finite-part reduction}
\label{gamma:sec}

The abstract reduction used in this section is the classical
functional-monomial decomposition of differential polynomials.
It appears explicitly in Guo--Regensburger--Rosenkranz
\cite[Proposition~4.10]{Guo2014}, with an earlier reduction rule in
Rosenkranz--Regensburger \cite[Lemma~5]{Rosenkranz2008}.
We give its short proof to specify the projection exactly.  The
special-function contribution is its application to the Gamma
hierarchy, the coefficient sectors in the next subsection, and the
endpoint and spectral formulas that follow.

\subsection{A canonical linear normal form}

Let $\mathscr R=\Q[u_0,u_1,u_2,\ldots]$, with derivation
$D u_r=u_{r+1}$ and weight $\wt(u_r)=r+1$.
The factor degree of $\prod u_r^{a_r}$ is $\sum a_r$.
A monomial is called functional if it is a power of $u_0$ (including
$1$), or its largest positive derivative index occurs at least twice.
Let $\mathscr F$ be their rational span.

\begin{theorem}[Classical normal form, with a fixed primitive]
\label{gamma:normal}
Every $P\in\mathscr R$ has a unique decomposition
\begin{equation}\label{gamma:decomp}
 P=\NN(P)+D\HH(P),\qquad \NN(P)\in\mathscr F,
 \qquad \HH(P)\text{ has zero constant term}.
\end{equation}
Both operators are rational and linear.  Reduction preserves factor
degree; $\NN$ preserves weight and $\HH$ lowers weight by one.
If $p(m)$ is the partition function, then
\begin{equation}\label{gamma:dimension}
 \dim_\Q\mathscr F_m=p(m)-p(m-1)\quad(m\ge2),
 \qquad \dim\mathscr F_0=\dim\mathscr F_1=1.
\end{equation}
\end{theorem}
\begin{proof}
If a monomial $M$ has a unique largest index $r\ge1$, replace that
factor $u_r$ by $u_{r-1}$, obtaining $\widetilde M$.
If $a$ is the multiplicity of $u_{r-1}$ in $\widetilde M$, then
\[
 M=\frac1aD\widetilde M+
      \left(M-\frac1aD\widetilde M\right).
\]
Every remaining monomial has largest derivative index at most $r-1$:
all differentiations of a factor $u_{r-1}$ give precisely the $a$
copies of $M$ that were removed.  Induction proves existence and
termination.  It also proves the grading assertions.

For uniqueness, a nonconstant polynomial $H$ of largest derivative
index $r$ has
$DH=u_{r+1}\partial H/\partial u_r+$ terms of order at most $r$.
Its terms of order $r+1$ have that largest index exactly once and
cannot belong to $\mathscr F$.  Hence $D\mathscr R\cap\mathscr F=0$,
and the kernel of $D$ consists of constants.  Fixing the constant of
$H$ proves uniqueness.  There are $p(m)$ monomials of weight $m$;
for $m\ge2$, $D$ is injective on the positive-weight space of weight
$m-1$.  This proves \eqref{gamma:dimension}.
\end{proof}

The quotient $\mathscr R/D\mathscr R$ in this theorem is a
\emph{vector-space quotient}.  For example $\NN(u_1)=0$, but
$\NN(u_1^2)=u_1^2$.  Thus $\NN$ is not a ring homomorphism.
No multiplication of congruences modulo $D\mathscr R$ is permitted
without a separate justification.

Now put $u_r=\psi^{(r)}(x)$.  At zero the convergent Laurent expansion is
\begin{equation}\label{gamma:local}
 \psi(x)=-\frac1x-\gamma+
       \sum_{j\ge1}(-1)^{j+1}\zeta(j+1)x^j,
       \qquad |x|<1.
\end{equation}
It follows from the standard log-Gamma series \cite[\S5.7]{DLMF};
all higher local expansions follow by differentiation.
Define the explicit endpoint functional
\begin{equation}\label{gamma:Delta}
 \Delta H=H\bigl(\psi(1),\psi'(1),\ldots\bigr)
             -[x^0]H\bigl(\psi(x),\psi'(x),\ldots\bigr).
\end{equation}
Then \eqref{scope:boundary} gives
\begin{equation}\label{gamma:fp-normal}
 \FP\int_0^1P(\psi,\psi',\ldots)\dd x
 =\FP\int_0^1\NN(P)(\psi,\psi',\ldots)\dd x+\Delta\HH(P).
\end{equation}
This is a finite algebraic recipe: for a monomial of weight $w$, write
it as $x^{-w}\prod_r(x^{r+1}\psi^{(r)}(x))^{a_r}$ and take its
coefficient of $x^w$.  Only terms through degree $w$ in these ordinary
series are needed.  The resulting endpoint belongs to
$\Q[\gamma,\zeta(2),\zeta(3),\ldots]$.

\subsection{The inverse-Gamma hierarchy and its balance identity}

Define
\begin{equation}\label{gamma:A}
 \A_m(x)=(-1)^m\Gamma(x)\frac{d^m}{dx^m}\frac1{\Gamma(x)},
 \qquad \A_0=1.
\end{equation}
As universal differential polynomials,
\begin{equation}\label{gamma:recurrence}
 \A_{m+1}=u_0\A_m-D\A_m,
 \qquad
 \sum_{m\ge0}\A_m\frac{t^m}{m!}
 =\exp\left(\sum_{r\ge0}\frac{(-1)^ru_rt^{r+1}}{(r+1)!}\right).
\end{equation}
The generating function is simply $\Gamma(x)/\Gamma(x-t)$.
Writing $B_j=\Gamma^{(j)}(x)/\Gamma(x)$ gives the exact balance
identity
\begin{equation}\label{gamma:balanced}
 \A_m=B_r\A_{m-r}
       -D\sum_{j=0}^{r-1}B_j\A_{m-1-j},\qquad 0\le r\le m.
\end{equation}
Indeed $DB_j=B_{j+1}-u_0B_j$ and
$D\A_k=u_0\A_k-\A_{k+1}$, so differentiating the sum telescopes.
Also $B_m(u)=(-1)^m\A_m(-u)$.  Taking $r=m$ in
\eqref{gamma:balanced} and then the normal form proves
\begin{equation}\label{gamma:parity}
 \NN(\A_m)(u)=(-1)^m\NN(\A_m)(-u).
\end{equation}
Only factor degrees congruent to $m$ modulo two survive.

The first reductions illustrate what must still be integrated:
\begin{align}
 \NN(\A_2)&=u_0^2,&
 \NN(\A_3)&=u_0^3,\notag\\
 \NN(\A_4)&=u_0^4-u_1^2,&
 \NN(\A_5)&=u_0^5-5u_0u_1^2,\notag\\
 \NN(\A_6)&=u_0^6-15u_0^2u_1^2+u_2^2,\label{gamma:firstforms}\\
 \NN(\A_7)&=u_0^7-35u_0^3u_1^2+7u_0u_2^2,\notag\\
 \NN(\A_8)&=u_0^8-70u_0^4u_1^2+28u_0^2u_2^2-7u_1^4-u_3^2.
 \notag
\end{align}
The full differential identities, including primitives, are exact
polynomial identities.  The accompanying code supplies them through
weight ten; the same finite algorithm applies at every weight.

\subsection{The inherited scalar bridge, with its normalization proved}

Let
\begin{equation}\label{gamma:c}
 \frac1{\Gamma(1+u)}=\sum_{j\ge0}c_j u^j,
 \qquad c_0=1,\quad c_1=\gamma.
\end{equation}
The source reports establish the all-depth Gamma-quotient bridge
\cite[\S3]{SourceCanonical}.  We reproduce a proof here because its
Laurent sign and scalar normalization enter the sixth moment.
Define
\begin{equation}\label{gamma:C}
 \mathcal C_m(x)=m![t^m]
       \frac{\Gamma(1-t)\Gamma(x)}{\Gamma(x-t)}.
\end{equation}
For $m\ge2$ it extends continuously to $[0,1]$, and
\begin{equation}\label{gamma:Cintegral}
 \int_0^1\mathcal C_m(x)\dd x=m!(c_m-\lambda_{m-1}).
\end{equation}

Here is a direct Mellin normalization.  Set
$L(v)=-\log(1-e^{-v})$ and
\[
 V_d(v)=L(v)^d-\boldsymbol1_{(0,1)}(v)(-\log v)^d.
\]
The generating identity for strict unit-index harmonic sums and
integration by parts give, initially for $\re u>0$,
\begin{equation}\label{gamma:mellinheight}
 \zeta(1+u,\{1\}^{d-1})
 =\frac1{\Gamma(1+u)u^d}
  +\frac{u}{d!\Gamma(1+u)}\int_0^\infty v^{u-1}V_d(v)\dd v.
\end{equation}
For clarity, the unsplit integral is
$[\Gamma(u)d!]^{-1}\int_0^\infty v^{u-1}L(v)^d\dd v$.
It follows by using
\[
 \sum_{n\ge1}H_{n-1}(\{1\}^{d-1})e^{-nv}
 =\frac{e^{-v}L(v)^{d-1}}{(d-1)!(1-e^{-v})}
 =-\frac1{d!}\frac d{dv}L(v)^d.
\]
The subtracted integral in \eqref{gamma:mellinheight} is holomorphic
near $u=0$: near zero $V_d(v)=O(v(1+|\log v|)^{d-1})$, and at
infinity it decays exponentially.  Consequently
\begin{equation}\label{gamma:lambda-integral}
 \lambda_d=c_{d+1}+\frac1{d!}\int_0^\infty\frac{V_d(v)}v\dd v.
\end{equation}

For the other side put
$F(x,t)=\Gamma(1+t)\Gamma(x)/\Gamma(x+t)$ and
\[
 \Phi(t)=\int_0^1\{F(x,t)-1+t(\psi(x)+\gamma)\}\dd x.
\]
The residue $t/x$ of $F$ cancels the residue $-t/x$ of the last
term.  The remaining integrand is jointly holomorphic in $t$ near
zero and continuous in $x\in[0,1]$.
Euler's beta integral, combined with the digamma integral before
interchanging integrations, yields
\begin{equation}\label{gamma:Phi}
 \Phi(t)=t\int_0^\infty
 \frac{(1-e^{-v})^t-1-\boldsymbol1_{(0,1)}(v)(v^t-1)}v\dd v.
\end{equation}
To verify the subtraction constant, for $\re t>0$ integration in
$x$ first gives
\[
 t\int_0^\infty
 \left\{\frac{(1-e^{-v})^t-1}{v}
              +\frac{e^{-v}}{1-e^{-v}}\right\}\dd v-1.
\]
With lower cutoff $\delta$, the second integral is
$-\log(1-e^{-\delta})$ and
$\int_\delta^1(v^t-1)\dd v/v=(1-\delta^t)/t+\log\delta$.
Their limits give exactly \eqref{gamma:Phi}.
The latter integral continues holomorphically to a neighborhood of
zero by its displayed local subtraction.  Comparing the coefficient
of $t^m$ on its two sides gives
$\int\mathcal C_m=-m\int V_{m-1}(v)\dd v/v$, proving
\eqref{gamma:Cintegral}.  This proof also justifies all coefficient
extractions as ordinary integrals.

\begin{theorem}[All inverse-Gamma finite parts]
\label{gamma:scalar}
For $m\ge2$ put $\mathcal L_m=\FP\int_0^1\A_m(x)\dd x$.
Then
\begin{equation}\label{gamma:L}
 \boxed{\displaystyle
 \mathcal L_m=K_m-m!\sum_{d=1}^{m-1}
       (-1)^{m-1-d}c_{m-1-d}\lambda_d,\qquad
 K_{2n}=\frac{(-1)^n\pi^{2n}}{2n+1},\quad K_{2n+1}=0.}
\end{equation}
This is the scalar bridge already implicit in the incoming all-order
Gamma formulas; the normal-form and spectral applications below
are the extensions developed here.
\end{theorem}
\begin{proof}
The coefficientwise finite part of the normalized quotient is
\[
 1+\gamma t+\sum_{m\ge2}(c_m-\lambda_{m-1})t^m
 =\frac1{\Gamma(1+t)}-\sum_{d\ge1}\lambda_d t^{d+1}.
\]
Here $\FP\int\mathcal C_1=\gamma$, since
$\mathcal C_1=\psi+\gamma$ and $\FP\int\psi=0$.
Multiply by $1/\Gamma(1-t)=\sum(-1)^jc_jt^j$.
Gamma reflection turns the scalar product into
$\sin(\pi t)/(\pi t)$.  Extracting $m![t^m]$ proves
\eqref{gamma:L}.
\end{proof}

If
$\NN(\A_m)=u_0^m+R_m$ and $H_m=\HH(\A_m)$, the immediate
all-order moment identity is
\begin{equation}\label{gamma:allmoments}
 \boxed{\displaystyle
 Q_m:=\FP\int_0^1\psi(x)^m\dd x
   =\mathcal L_m-\Delta H_m
       -\FP\int_0^1R_m(\psi,\psi',\ldots)\dd x.}
\end{equation}
Every monomial of $R_m$ contains a derivative.  Section~\ref{spectral:sec}
will evaluate that remainder by convergent harmonic coordinates and
provide its explicit antiderivative, so \eqref{gamma:allmoments} is
part of a terminating exact procedure rather than an instruction to
leave mixed integrals unnamed.

% Fragment for the parent article.  Uses its differential-polynomial
% notation A_m, N, D and u_r; replace labels/macros as needed.
\subsection{Two exact coefficient sectors at every degree}
\label{nfsector:section}

Let $\mathcal A_m$ be defined by
\[
 \mathcal A_0=1,\qquad
 \mathcal A_{m+1}=u_0\mathcal A_m-D\mathcal A_m,
 \qquad D u_r=u_{r+1}.
\]
Equivalently its exponential generating series is
\begin{equation}\label{nfsector:Bell}
 \mathcal A(t)=\sum_{m\ge0}\mathcal A_m\frac{t^m}{m!}
 =\exp\left(\sum_{r\ge0}
       \frac{(-1)^r u_r t^{r+1}}{(r+1)!}\right).
\end{equation}
All series in this subsection are formal power series.  Therefore each
coefficient computation is finite.  Write $\mathcal N$ for the classical
functional-monomial projection along the vector space $D\mathbb Q[u_0,u_1,\ldots]$.
It retains powers of $u_0$ and monomials in which the largest positive
derivative index occurs at least twice.  This projection is linear and
need not respect products.

\begin{lemma}[Translation and centered generation]\label{nfsector:Appell}
For an independent constant $c$,
\begin{equation}\label{nfsector:Appellformula}
 \mathcal N(\mathcal A_m)(u_0+c,u_1,\ldots)
 =\sum_{j=0}^m\binom mj c^{m-j}\mathcal N(\mathcal A_j)(u_0,u_1,\ldots).
\end{equation}
Also
\begin{equation}\label{nfsector:centered}
 \mathcal N\mathcal A(t)
 =\mathcal N\exp\left(\sum_{k\ge0}
       \frac{u_{2k}t^{2k+1}}{2^{2k}(2k+1)!}\right).
\end{equation}
\end{lemma}
\begin{proof}
Substitution $u_0\mapsto u_0+c$ commutes with $D$, and preserves the
span of the functional monomials.  Uniqueness of the direct-sum
decomposition therefore implies that it commutes with $\mathcal N$.
Formula~\eqref{nfsector:Bell} acquires the factor $e^{ct}$ under this
substitution, proving \eqref{nfsector:Appellformula}.

For the second assertion, apply the formal translation $e^{tD/2}$
to \eqref{nfsector:Bell}.  Directly summing its exponent gives
\[
 e^{tD/2}\mathcal A(t)
 =\exp\left(\sum_{k\ge0}
       \frac{u_{2k}t^{2k+1}}{2^{2k}(2k+1)!}\right).
\]
Every coefficient of $(e^{tD/2}-1)\mathcal A(t)$ belongs to the image
of $D$, proving the claim.  In particular every surviving monomial
has factor degree congruent to its weight modulo two: the centered
exponent contains only variables of odd weight.
\end{proof}

\begin{theorem}[All polygamma-square coefficients]\label{nfsector:squares}
For every $r\ge1$ and $m\ge2r+2$,
\begin{equation}\label{nfsector:squarecoef}
 [u_0^{m-2r-2}u_r^2]\,\mathcal N(\mathcal A_m)
       =(-1)^r\binom{m}{2r+2}.
\end{equation}
Thus every square term with exactly two factors of positive derivative
order is known at all weights.
\end{theorem}
\begin{proof}
The operator $D$ and the reduction preserve factor degree.  At weight
$2r+2$, the quadratic part of $\mathcal A_{2r+2}$ is
\[
 \frac{(2r+2)!}{2}
 \sum_{a+b=2r}\frac{u_a u_b}{(a+1)!(b+1)!}.
\]
Repeated integration by parts gives
$\mathcal N(u_a u_{2r-a})=(-1)^{a-r}u_r^2$.
The coefficient in question is consequently
\[
 \frac12\sum_{j=1}^{2r+1}
       (-1)^{j-1-r}\binom{2r+2}{j}=(-1)^r,
\]
because the full alternating binomial sum is zero and its two
endpoint terms are both one.  Formula~\eqref{nfsector:Appellformula},
with the original $u_0$ set to zero, now supplies the factor
$\binom{m}{2r+2}$.
\end{proof}

\begin{theorem}[The complete digamma--trigamma sector]
\label{nfsector:trigamma}
After the normal form has been computed, set $u_r=0$ for every $r\ge2$.
Then, for all $m\ge0$,
\begin{equation}\label{nfsector:sector}
 \left.\mathcal N(\mathcal A_m)\right|_{u_2=u_3=\cdots=0}
 =u_0^m-
 \sum_{s=1}^{\lfloor m/4\rfloor}
 \binom{m}{4s}\frac{(4s-1)!!}{4s^2-1}
       u_0^{m-4s}u_1^{2s}.
\end{equation}
Equivalently, put $y=u_1t^2/2$ and define the entire function
$\operatorname{Shi}(y)=\int_0^y\sinh(v)\,dv/v$.  The exponential
generating series of \eqref{nfsector:sector} is
\begin{equation}\label{nfsector:Shi}
 \left.\sum_{m\ge0}\mathcal N(\mathcal A_m)\frac{t^m}{m!}
       \right|_{u_2=u_3=\cdots=0}
 =\frac{e^{u_0t}}2
       \left(\cosh y+\frac{\sinh y}{y}-y\operatorname{Shi}(y)\right),
\end{equation}
where $\sinh y/y$ has its removable value one at zero.
\end{theorem}
\begin{proof}
Let $I$ be the monomial ideal generated by $u_r u_s$ for $r,s\ge2$.
It is stable under $D$.  The stronger property needed here is
\begin{equation}\label{nfsector:ideal}
                   \mathcal N(I)\subseteq I.
\end{equation}
Indeed, an irreducible monomial in $I$ is retained.  If a reducible
monomial in $I$ has a unique largest index $r$, then $r\ge3$.
Replacing $u_r$ by $u_{r-1}$ in its reduction primitive preserves at
least two factors of index at least two.  The primitive and its
derivative therefore remain in $I$.  Induction through the reduction
proves \eqref{nfsector:ideal}.  This argument is essential; stability
under $D$ alone would not imply stability under an arbitrary projection.

Modulo $I$, the Bell series contains at most one factor of derivative
order at least two:
\begin{equation}\label{nfsector:modI}
 \mathcal A(t)\equiv
 e^{u_0t-u_1t^2/2}
 \left(1+\sum_{r\ge2}\frac{(-1)^ru_rt^{r+1}}{(r+1)!}\right)
 \pmod I.
\end{equation}
For integers $l,k\ge0$ and $r\ge2$, integration by parts yields,
modulo the vector subspace $D\mathbb Q[u_0,u_1,\ldots]+I$,
\begin{equation}\label{nfsector:onehigh}
 u_0^l u_1^k u_r\equiv
 \frac{(-1)^{r-1}(l)_{\underline{r-1}}}{k+r-1}
       u_0^{l-r+1}u_1^{k+r}
       \quad(l\ge r-1).
\end{equation}
For $l<r-1$ the class is zero.  Here the underline denotes a falling
factorial.  To prove the formula, for $r\ge3$ differentiate
$u_0^l u_1^k u_{r-1}$: the term containing $u_2u_{r-1}$ belongs
to $I$, leaving a factor $-l$ and reducing $r$ by one.
The final step is
\[
 u_0^l u_1^k u_2\equiv
       -\frac{l}{k+1}u_0^{l-1}u_1^{k+2}.
\]
If there are too few factors of $u_0$, a zero coefficient appears
before this last step.

Let $d_M=[u_1^M]\mathcal N(\mathcal A_{2M})$ for $M\ge2$.
The term without a higher derivative in \eqref{nfsector:modI}
contributes $(2M)!(-1)^M/(2^MM!)$.  For a term with one $u_r$,
the necessary indices are $l=r-1$ and $k=M-r$, where $2\le r\le M$.
Formula~\eqref{nfsector:onehigh} therefore gives
\begin{align*}
 \frac{d_M}{(2M)!}
 &=\frac{(-1)^M}{2^MM!}
   +\frac{(-1)^M}{M-1}
       \sum_{r=2}^{M}
       \frac{(-1)^{r-1}}{(M-r)!2^{M-r}(r+1)!}\\
 &=\frac{(-1)^M\bigl((-1)^{M+1}-1\bigr)}
       {2^{M+1}(M-1)(M+1)!}.
\end{align*}
For the last equality, use the binomial expansion of
$(1/2-1)^{M+1}$ and remove its terms of indices zero, one and two.
Thus $d_M=0$ for odd $M$, while
\[
 d_M=-\frac{(2M-1)!!}{M^2-1}\qquad(M\ge2\text{ even}).
\]
The pure $u_0$ coefficient is one; there is no retained linear $u_1$
term.  The Appell identity now proves \eqref{nfsector:sector}.

Finally, its generating factor is
\[
 1-\sum_{s\ge1}\frac{y^{2s}}{(2s)!((2s)^2-1)}.
\]
Expanding $\cosh y$, $\sinh y/y$ and $y\operatorname{Shi}(y)$
shows that this is the bracket in \eqref{nfsector:Shi} divided by two.
This also proves that the displayed apparent singularity is removable.
\end{proof}

The order of operations in \eqref{nfsector:sector} is part of its
statement: normal reduction precedes specialization.  No assertion is
being made that the normal projection is multiplicative or that one
may impose $u_2=u_3=\cdots=0$ before integration by parts.
On substituting $u_r=\psi^{(r)}(x)$, these formulas give every
digamma--trigamma coefficient and every pure polygamma-square coefficient
in the finite-part moment reductions.  All omitted normal-form terms
in \eqref{nfsector:sector} contain at least two factors of derivative
order at least two.  Their integrals and the exact endpoint functional
must still be retained in the full identity.

% Independent verification note for the parent arithmetic reduction.
\section[Hurwitz products, spectral jets, and Stieltjes derivative correlations]{Hurwitz products, spectral jets,\\and Stieltjes derivative correlations}
\label{hprod:sec}

For integers $k,l\ge2$, put $W=k+l$ and
$C_{kl}=\binom{W-2}{k-1}$.  In the original coordinate $x$, define
\[
 I_{kl}=\operatorname{FP}_{\epsilon\downarrow0}
 \int_\epsilon^1\zeta(k,x)\zeta(l,x)\,dx.
\]
Here the finite part retains the constant coefficient after all negative
powers of $\epsilon$ and all logarithmic divergences are removed.
Then
\[
 I_{kl}=
 \begin{cases}
 (-1)^kC_{kl}\bigl\{[2H_{W-2}-H_{k-1}-H_{l-1}]\zeta(W-1)
                    +2\zeta'(W-1)\bigr\},&k\equiv l\pmod2,\\[3pt]
 (-1)^kC_{kl}[H_{l-1}-H_{k-1}]\zeta(W-1),&k\not\equiv l\pmod2.
 \end{cases}
\]
In particular
\[
 I_{22}=2\zeta(3)+4\zeta'(3),\qquad
 I_{33}=-7\zeta(5)-12\zeta'(5),\qquad
 I_{23}=\tfrac32\zeta(4),\qquad
 I_{24}=\tfrac{16}{3}\zeta(5)+8\zeta'(5).
\]
For polygamma products one multiplies $I_{kl}$ by
$(-1)^{k+l}(k-1)!(l-1)!$.

\paragraph{Proof and normalization.}
The classical product kernel \cite[Remark~1]{Shpot2016}, also recorded in \cite[Chapter~7]{SourceManuscript}, underlies the following proof.
The classical Mikol\'as/Hurwitz product integral, already reproduced
in the ProveIt manuscript's coincident-endpoint chapter, continues
meromorphically as
\begin{align*}
 F(s,t)&=2\Gamma(1-s)\Gamma(1-t)(2\pi)^{s+t-2}
 \cos\!\left(\frac{\pi(s-t)}2\right)\zeta(2-s-t)\\
 &=\frac{\pi\Gamma(s+t-1)\zeta(s+t-1)}{\Gamma(s)\Gamma(t)}
   \{\csc(\pi s)+\csc(\pi t)\}.
\end{align*}
The second expression follows from the Riemann functional equation and
Gamma reflection.  To relate its Laurent coefficient to the coordinate
finite part, write
\[
 \zeta(s,x)=x^{-s}+\sum_{j=0}^{J}c_j(s)x^j+O(x^{J+1}),
 \qquad c_j(s)=\frac{(-1)^j(s)_j}{j!}\zeta(s+j),
\]
uniformly for $s$ in a small neighborhood of $k$, with $J$ sufficiently
large.  There is an analogous expansion near $l$.  Uniformly subtract
the resulting local terms before integrating.  Along
$s=k+u,t=l+u$, the only local terms whose exponent crosses $-1$ are
\[
 c_{k-1}(l+u)x^{-1-u},\qquad
 c_{l-1}(k+u)x^{-1-u}.
\]
Their analytic integrals are respectively
$-c_{k-1}(l+u)/u$ and $-c_{l-1}(k+u)/u$, whereas their coordinate
finite parts at $u=0$ are zero.  All other subtracted monomials have
denominators nonzero at zero and commute with evaluation.  Therefore
\[
 I_{kl}=[u^0]F(k+u,l+u)
       +c'_{k-1}(l)+c'_{l-1}(k).
\]
This is the precise local correction; extracting a spectral Laurent
constant alone would give the wrong coordinate finite part.

If $k,l$ have the same parity, let $\sigma=(-1)^k=(-1)^l$.
Expansion of the cosecant representation gives
\[
 [u^0]F(k+u,l+u)
 =2\sigma C_{kl}\{[2\psi(W-1)-\psi(k)-\psi(l)]\zeta(W-1)
                   +2\zeta'(W-1)\}.
\]
The two coefficient derivatives have sum equal to the negative of one
half of this expression.  If $k,l$ have opposite parity, the two
cosecants cancel identically on the chosen line, so its Laurent
constant is zero; the coefficient derivatives give
$(-1)^kC_{kl}[\psi(l)-\psi(k)]\zeta(W-1)$.  Finally
$\psi(n)=H_{n-1}-\gamma$ proves the displayed formulas.

\subsection{All order derivatives from a holomorphic local generator}

The coordinate correction extends to every independent order derivative.
Continue to write
$c_j(s)=(-1)^j(s)_j\zeta(s+j)/j!$, and define
\begin{equation}\label{hprod:completed-generator}
 \mathscr H_{kl}(s,t)=F(s,t)
 +\frac{c_{k-1}(t)}{s-k}+\frac{c_{l-1}(s)}{t-l}.
\end{equation}
The singularities displayed separately in this expression are removable
in their sum near $(k,l)$.

\begin{theorem}[Every independent order jet]
\label{hprod:all-jets}
For every $k,l\ge2$, $\mathscr H_{kl}$ is holomorphic in a
neighborhood of $(k,l)$.  For all nonnegative integers $a,b$,
\begin{equation}\label{hprod:jet-fp}
 \left.\partial_s^a\partial_t^b\mathscr H_{kl}(s,t)
 \right|_{s=k,t=l}
 =\operatorname{FP}\int_0^1
 \left.\partial_s^a\zeta(s,x)\right|_{s=k}
 \left.\partial_t^b\zeta(t,x)\right|_{t=l}\,dx.
\end{equation}
The finite part is taken in the same original coordinate $x$ for all
orders.  Every jet on the left can be obtained by finite Taylor
arithmetic in Gamma functions and the ordinary Riemann zeta function.
\end{theorem}

\begin{proof}
Choose small neighborhoods of $k$ and $l$ avoiding the spectral pole at
$1$, and subtract a sufficiently
long Taylor polynomial of each $g_s(x)=\zeta(s,1+x)$ at $x=0$.
The product consists of an integrable holomorphic remainder, the
monomial $x^{-s-t}$, and finitely many cross monomials
$c_j(t)x^{j-s}$ and $c_j(s)x^{j-t}$.  Near $(k,l)$, all
denominators obtained by integrating these monomials are nonzero
except $k-s$ in the first cross family and $l-t$ in the second.
Those two integrals are exactly
$-c_{k-1}(t)/(s-k)$ and $-c_{l-1}(s)/(t-l)$; adding them
back proves holomorphy of \eqref{hprod:completed-generator}.

Uniform remainder estimates permit every fixed number of parameter
derivatives under the integrable remainder: the resulting bounds have
the form $C x^{-\delta}(1+|\log x|^N)$ with $0<\delta<1$.
For each nonresonant monomial the coordinate finite part and parameter
differentiation agree, since its denominator is nonzero.  All derivatives
of either removed resonant monomial, evaluated at $(k,l)$, are
constant multiples of $x^{-1}\log^j x$.  Their integrals over
$[\epsilon,1]$ have constant coefficient zero for every $j\ge0$.
Thus they contribute zero to the fixed-coordinate finite part at every
jet order, proving \eqref{hprod:jet-fp}.
\end{proof}

There is a useful version of \eqref{hprod:completed-generator} with
its removable singularities already resolved.  Put
\[
 P(s,t)=\frac{\Gamma(s+t-1)}{\Gamma(s)\Gamma(t)}\zeta(s+t-1),
 \quad E(u)=\frac{\pi u}{\sin(\pi u)},\quad
 \sigma_k=(-1)^k,
\]
and write $s=k+u$, $t=l+v$.  Since
$c_{k-1}(t)=-\sigma_kP(k,t)$, one obtains
\begin{equation}\label{hprod:divided-differences}
 \begin{split}
 \mathscr H_{kl}(k+u,l+v)
 ={}&\sigma_k\frac{P(k+u,l+v)E(u)-P(k,l+v)}u\\
 &+\sigma_l\frac{P(k+u,l+v)E(v)-P(k+u,l)}v.
 \end{split}
\end{equation}
Both displayed quotients are holomorphic divided differences.
If
\[
 P(k+u,l+v)=\sum_{a,b\ge0}P_{ab}u^av^b,
 \qquad E(u)=\sum_{r\ge0}e_{2r}u^{2r},
\]
then $e_0=1$, $e_2=\pi^2/6$, $e_4=7\pi^4/360$, and
the jet in \eqref{hprod:jet-fp} is the finite expression
\begin{equation}\label{hprod:finite-jet-algorithm}
 a!b!\left\{
 \sigma_k\sum_{r=0}^{\lfloor(a+1)/2\rfloor}e_{2r}P_{a+1-2r,b}
 +\sigma_l\sum_{r=0}^{\lfloor(b+1)/2\rfloor}e_{2r}P_{a,b+1-2r}
 \right\}.
\end{equation}
Only Gamma/zeta Taylor coefficients through total degree $a+b+1$
are needed.  This gives an explicit terminating procedure rather than
an instruction to evaluate a difference of poles numerically.

\subsection{A first derivative identity with every correction}

Let
\[
 A_k=H_{2k-2}-H_{k-1},\qquad
 H_n^{(2)}=\sum_{j=1}^n\frac1{j^2}.
\]
The first diagonal order jet is
\begin{equation}\label{hprod:first-order-jet}
 \begin{split}
 &\operatorname{FP}\int_0^1\zeta(k,x)
     \left.\partial_s\zeta(s,x)\right|_{s=k}\,dx\\
 &\quad=\frac{(-1)^k}{2}\binom{2k-2}{k-1}
 \biggl\{
 \left[3A_k^2+4\zeta(2)-3H_{2k-2}^{(2)}+H_{k-1}^{(2)}\right]
        \zeta(2k-1)\\
 &\hspace{46mm}+6A_k\zeta'(2k-1)+3\zeta''(2k-1)
 \biggr\}.
 \end{split}
\end{equation}
In particular,
\begin{equation}\label{hprod:trigamma-order-jet}
 \operatorname{FP}\int_0^1\zeta(2,x)
     \left.\partial_s\zeta(s,x)\right|_{s=2}\,dx
 =\bigl(4\zeta(2)-2\bigr)\zeta(3)
   +3\zeta'(3)+3\zeta''(3).
\end{equation}

\begin{proof}
Set
\[
 B(s)=\frac{\Gamma(2s-1)}{\Gamma(s)^2}\zeta(2s-1),\qquad
 D(s)=\frac{\Gamma(s+k-1)}{\Gamma(s)\Gamma(k)}\zeta(s+k-1).
\]
Along the diagonal through $(k,k)$, the functional equation and
\eqref{hprod:completed-generator} give
\[
 \mathscr H_{kk}(s,s)=2\pi\csc(\pi s)B(s)
                   -\frac{2(-1)^kD(s)}{s-k}.
\]
Since $B(k)=D(k)$, the coefficient of $s-k$ is
\[
 (-1)^k\left\{B''(k)-D''(k)+\frac{\pi^2}{3}B(k)\right\}.
\]
Symmetry of $\mathscr H_{kk}$ makes the desired single order
derivative one half of that coefficient.  Logarithmic differentiation
of the two Gamma quotients yields
\begin{align*}
 \frac{B''(k)}{C_{kk}}
 &=\left[4A_k^2+2\zeta(2)-4H_{2k-2}^{(2)}+2H_{k-1}^{(2)}\right]
       \zeta(2k-1)\\
 &\qquad+8A_k\zeta'(2k-1)+4\zeta''(2k-1),\\
 \frac{D''(k)}{C_{kk}}
 &=\left[A_k^2-H_{2k-2}^{(2)}+H_{k-1}^{(2)}\right]
       \zeta(2k-1)
       +2A_k\zeta'(2k-1)+\zeta''(2k-1).
\end{align*}
Subtraction and $\pi^2/3=2\zeta(2)$ prove
\eqref{hprod:first-order-jet}; the $k=2$ specialization gives
\eqref{hprod:trigamma-order-jet}.
\end{proof}

\subsection{Products of ordinary derivatives of generalized Stieltjes constants}

The completed kernel also evaluates all quadratic products of positive
ordinary derivative orders of generalized Stieltjes constants.  Use the
standard convention
\[
 \zeta(1+u,x)=\frac1u+
 \sum_{n=0}^{\infty}\frac{(-1)^n}{n!}\gamma_n(x)u^n.
\]
For integers $r\ge1$ and $0\le j\le r$, put
\[
 p_{r,j}=\left.\frac{d^j}{du^j}(1+u)_r\right|_{u=0},
\]
and define the constant-coefficient differential operator
\begin{equation}\label{hprod:stieltjes-operator}
 T_{n,r}^{(s)}=
 \sum_{j=0}^{\min(n,r)}\binom nj p_{r,j}\partial_s^{n-j}.
\end{equation}
Its coefficients are explicit positive integers; in particular
$p_{r,0}=r!$.  They may equally be computed from
$p_{r,j}=j!\,[u^j](1+u)_r$.

\begin{corollary}[Every quadratic derivative-Stieltjes finite part]
\label{hprod:stieltjes-products}
For all Stieltjes indices $m,n\ge0$ and positive ordinary derivative
orders $r,t\ge1$,
\begin{equation}\label{hprod:stieltjes-family}
 \begin{split}
 &\operatorname{FP}\int_0^1
       \gamma_m^{(r)}(x)\gamma_n^{(t)}(x)\,dx\\
 &\quad=(-1)^{m+n+r+t}
 \left.T_{m,r}^{(s)}T_{n,t}^{(v)}
       \mathscr H_{r+1,t+1}(s,v)\right|_{s=r+1,v=t+1}.
 \end{split}
\end{equation}
The right side is evaluated by the finite algorithm
\eqref{hprod:finite-jet-algorithm}.  Superscripts $(r)$ and $(t)$
on $\gamma$ mean ordinary derivatives with respect to its argument
$x$, not derivatives with respect to its Stieltjes index.
\end{corollary}

\begin{proof}
For $r\ge1$, differentiation with respect to $x$ eliminates the
spectral pole $1/u$.  The Hurwitz differentiation formula gives
\[
 \sum_{n=0}^{\infty}\frac{(-1)^n}{n!}
       \gamma_n^{(r)}(x)u^n
 =(-1)^r(1+u)_r\zeta(r+1+u,x).
\]
Taking the coefficient of $u^n$ and applying the product rule yields
\begin{equation}\label{hprod:stieltjes-to-jets}
 \gamma_n^{(r)}(x)=(-1)^{n+r}
 \left.T_{n,r}^{(s)}\zeta(s,x)\right|_{s=r+1}.
\end{equation}
Multiply the two finite expressions
\eqref{hprod:stieltjes-to-jets}, use linearity of the specified
coordinate finite part, and apply Theorem~\ref{hprod:all-jets}
term by term.  This proves \eqref{hprod:stieltjes-family}.
\end{proof}

For the first nontrivial mixed-index example,
\[
 \gamma_0'(x)=-\zeta(2,x),\qquad
 \gamma_1'(x)=\zeta(2,x)+
       \left.\partial_s\zeta(s,x)\right|_{s=2}.
\]
Combining $I_{22}$ with \eqref{hprod:trigamma-order-jet} gives
\begin{equation}\label{hprod:stieltjes-first-example}
 \boxed{\operatorname{FP}\int_0^1
         \gamma_0'(x)\gamma_1'(x)\,dx
 =-4\zeta(2)\zeta(3)-7\zeta'(3)-3\zeta''(3).}
\end{equation}
This family is a consequence of the completed classical product kernel
and the standard Hurwitz derivative formula.  Its range here is
explicitly $r,t\ge1$; undifferentiated Stieltjes factors require the
separate completion at spectral order $1$ already developed in the
source manuscript.  No broad historical priority claim is attached to
the kernel or this coefficient-extraction method.

\paragraph{Independent checks.}
The local subtraction above also produces an ordinary integrable
function.  For direct numerical evaluation, put
$g_m(x)=\zeta(m,1+x)$ and
$R_{m,r}(x)=g_m(x)-\sum_{j=0}^{r-1}c_j(m)x^j$.  Then
\begin{align*}
 I_{kl}={}&-\frac1{W-1}
 +\sum_{j=0}^{k-2}\frac{c_j(l)}{j-k+1}
 +\sum_{j=0}^{l-2}\frac{c_j(k)}{j-l+1}\\
 &+\int_0^1\{g_k(x)g_l(x)+x^{-k}R_{l,k}(x)
                                +x^{-l}R_{k,l}(x)\}\,dx.
\end{align*}
Both apparent singularities in the last integrand are removable.  This
representation gives a numerical left side independent of the
Mikol\'as identity.  For $k=l=2$ it gives
$1.6116088347766411574767490370097567940\ldots$, and for
$k=l=3$ it gives
$-6.9156089198900340833548806001936246246\ldots$.

\paragraph{Source scope.}
The classical product integral should be credited to the source
literature; see, for example, Shpot--Chaudhary--Paris,
\emph{Integrals of products of Hurwitz zeta functions and the Casimir
effect in $\phi^4$ field theories}, arXiv:1603.00722v2, Remark~1,
and its cited Mikol\'as sources.  That paper's Theorem~2 concerns the
different ordinary integral of $\zeta(k,1+x)^2$.  The finite-part
identity here is an explicit integer-resonance consequence with its
coordinate correction proved.  No claim of exhaustive historical
novelty is made.

\section{Spectral coordinates for derivative-containing polynomials}
\label{spectral:sec}

The preceding normal forms reduce higher digamma moments to differential
polynomials containing at least one positive derivative.  For this
entire class, the following cancellation removes the convergence
subtractions usually needed in a Mittag--Leffler expansion.  All finite
parts in this section use the original coordinate $x$ at $x=0$.

Put $\mathcal R=\Q[u_0,u_1,\ldots]$ and
$\mathcal J=(u_1,u_2,\ldots)\subset\mathcal R$.  For $P\in\mathcal J$,
write
\[
 f_P(z)=P\bigl(\psi(z),\psi'(z),\ldots\bigr),\qquad
 f_P(-n+v)=\sum_{j=1}^{J}\frac{p_j(n)}{v^j}+O(1),
 \qquad r(n)=p_1(n).
\]
Here $J$ is any common upper bound for the pole orders; zero coefficients
are understood.  Let $\Delta a(n)=a(n)-a(n-1)$ for $n\geq1$.

\begin{theorem}[Unsubtracted expansion and finite-part coordinates]
\label{spectral:ML}
For every $P\in\mathcal J$,
\begin{equation}\label{spectral:MLseries}
 \boxed{\displaystyle
 f_P(z)=\sum_{n=0}^{\infty}\sum_{j=1}^{J}
                    \frac{p_j(n)}{(z+n)^j}.}
\end{equation}
The series is absolutely and locally uniformly convergent away from
$0,-1,-2,\ldots$.  For some fixed $A\geq0$, its coefficients satisfy
\begin{align}
 p_j(n)&=O\bigl(\log^A(n+2)\bigr),&
 \Delta p_j(n)&=O\bigl(n^{-1}\log^A(n+2)\bigr),\label{spectral:pest}\\
 r(n)&=O\bigl(n^{-1}\log^A(n+2)\bigr),&
 \Delta r(n)&=O\bigl(n^{-2}\log^A(n+2)\bigr).
 \label{spectral:rest}
\end{align}
Consequently the Dirichlet series
\begin{equation}\label{spectral:DP}
 D_P(s)=\sum_{n=1}^{\infty}\frac{\Delta r(n)}{n^s}
\end{equation}
is holomorphic for $\re s>-1$, and
\begin{equation}\label{spectral:FP}
 \boxed{\displaystyle
 \FP\int_0^1 f_P(x)\dd x
 =\sum_{n=1}^{\infty}\sum_{j=2}^{J}
       \frac{\Delta p_j(n)}{(j-1)n^{j-1}}+D_P'(0).}
\end{equation}
Every series in \eqref{spectral:FP}, including
$D_P'(0)=-\sum_{n\geq1}\Delta r(n)\log n$, converges absolutely.
\end{theorem}

\begin{proof}
The reflection formula and polygamma asymptotic expansions used below
are standard Gamma-function identities; see \cite[\S\S5.11,5.15]{DLMF}.
Set
\[
 K(v)=-\pi\cot(\pi v),\qquad K'(v)=K(v)^2+\pi^2.
\]
For any real $L$, every $K^{(a)}$, $a\geq1$, is a polynomial in $K$
divisible by $K^2+\pi^2$.  Since each monomial of $P$ contains a positive
derivative, there is a polynomial $T_L$ such that
\begin{equation}\label{spectral:frozen}
 P\bigl(L+K(v),K'(v),K''(v),\ldots\bigr)
       =\frac{\mathrm d}{\mathrm dv}T_L(K(v)).
\end{equation}
Indeed, divide the left side by $K^2+\pi^2$, integrate the resulting
polynomial in $K$, and use the displayed differential equation for $K$.
Thus the residue of \eqref{spectral:frozen} at $v=0$ is identically zero.

Fix $0<\rho<1/2$.  At the pole $-n$, reflection gives
\[
 \psi(-n+v)=K(v)+\log n+\eta_n(v),\qquad
 \eta_n(v)=\psi(n+1-v)-\log n.
\]
Extend the subscript to real $t\geq1$ in the rightmost formula.  Uniformly
on $|v|\leq\rho$, any fixed finite collection of $v$ derivatives satisfies
\[
 \partial_v^a\eta_t(v)=O(t^{-1}),\qquad
 \partial_t\partial_v^a\eta_t(v)=O(t^{-2})\quad(a\geq0).
\]
The first bound for $a=0$ includes the subtraction of $\log t$; the
second includes its derivative.  Coefficient extraction on $|v|=\rho$
therefore bounds every principal coefficient by a power of $\log t$,
and its $t$ derivative by $O(t^{-1}\log^A(t+2))$.  For the residue,
subtract \eqref{spectral:frozen} with $L=\log t$ before extracting the
coefficient.  Its residue is zero, while the difference is
$O(t^{-1}\log^A(t+2))$ and its $t$ derivative is
$O(t^{-2}\log^A(t+2))$.  The mean value theorem on $[n-1,n]$ gives
\eqref{spectral:pest}--\eqref{spectral:rest}; finitely many small $n$
do not affect the estimates.

For completeness, these local calculations leave no undetermined entire
function in \eqref{spectral:MLseries}.  Use the circles
$|w|=R_N=N+1/2$.  They stay at least $1/2$ from every integer.  In the
right half-plane the usual asymptotics apply directly; in the left
half-plane apply reflection to $1-w$.  On these circles one obtains
\[
 |\psi(w)|\leq C\log(R_N+2),\qquad
 |\psi^{(a)}(w)|\leq C_a
       \bigl(R_N^{-1}+e^{-c|\im w|}\bigr)\quad(a\geq1).
\]
The exponential term follows from the positive derivatives of cotangent;
pole avoidance makes the estimate uniform also when $|\im w|$ is small.
Hence
\[
 |f_P(w)|\leq C\log^A(R_N+2)
                  \bigl(R_N^{-1}+e^{-c|\im w|}\bigr).
\]
For fixed $z$,
\[
 \oint_{|w|=R_N}\frac{f_P(w)}{w-z}\dd w
     =O\!\left(\frac{\log^A(R_N+2)}{R_N}\right)\longrightarrow0,
\]
because $\int_0^{2\pi}e^{-cR_N|\sin\theta|}\dd\theta=O(R_N^{-1})$.
The residue at $w=-n$ of the integrand is
$-\sum_jp_j(n)/(z+n)^j$, and its residue at $w=z$ is $f_P(z)$.
The residue theorem gives \eqref{spectral:MLseries}.  Its absolute normal
convergence follows from \eqref{spectral:pest} and the extra factor
$n^{-1}$ in the estimate for $r(n)$.

Remove the finitely many principal terms at $0$.  The remaining series
can be integrated uniformly on $[0,1]$.  For $j\geq2$, the pole at zero
contributes $-p_j(0)/(j-1)$ to the finite part, and the other poles
contribute
\[
 \sum_{n\geq1}\frac{p_j(n)}{j-1}
                  \bigl(n^{1-j}-(n+1)^{1-j}\bigr).
\]
Summation by parts, with the vanishing boundary terms supplied by
\eqref{spectral:pest}, gives the higher-pole sum in
\eqref{spectral:FP}.  The simple pole at zero has finite part zero.
The remaining simple-pole terms give
\[
 \sum_{n\geq1}r(n)\log\frac{n+1}{n}
     =-\sum_{n\geq1}\Delta r(n)\log n=D_P'(0).
\]
Here the boundary term $r(N)\log(N+1)$ tends to zero.  The improved
difference estimate proves the asserted holomorphy and all absolute
convergence claims.
\end{proof}

\begin{corollary}[Anchored antiderivative]\label{spectral:primitive}
Let $\Omega=\C\setminus(-\infty,0]$ and use the principal logarithm.
For $z,z_0\in\Omega$, the function
\begin{align}\label{spectral:primitiveformula}
 F_P(z;z_0)={}&\sum_{n=0}^{\infty}r(n)
       \bigl(\operatorname{Log}(z+n)-\operatorname{Log}(z_0+n)\bigr)
       \notag\\
 &+\sum_{n=0}^{\infty}\sum_{j=2}^{J}\frac{p_j(n)}{j-1}
       \bigl((z_0+n)^{1-j}-(z+n)^{1-j}\bigr)
\end{align}
is locally normally convergent, satisfies $F_P(z_0;z_0)=0$, and has
$\partial_zF_P(z;z_0)=f_P(z)$.  Its difference of logarithms must be
retained for general complex anchors; it agrees with the principal
logarithm of the ratio when $z_0$ is positive real.
\end{corollary}

\begin{proof}
On a fixed compact subset of $\Omega\times\Omega$, the logarithm
difference is $O(n^{-1})$, while the inverse-power difference in the
$j$th term is $O(n^{-j})$.  The preceding coefficient bounds give normal
convergence.  Termwise differentiation and
\eqref{spectral:MLseries} finish the proof.
\end{proof}

\subsection{A finite conversion to convergent multiple-zeta coordinates}

Use the descending-index convention
\[
 \zeta(s_1,\ldots,s_\ell)
    =\sum_{n_1>\cdots>n_\ell\geq1}
       n_1^{-s_1}\cdots n_\ell^{-s_\ell}.
\]
The depth is $\ell$.  Write $H_N(\alpha_1,\ldots,\alpha_\ell)$ for the
same nested sum restricted to $N\geq n_1>\cdots>n_\ell\geq1$, and put
$H_N(\varnothing)=1$.  A coordinate derivative differentiates one
of the displayed arguments.  For the statements below, coefficients
belong to $\Q[\gamma,\zeta(2),\zeta(3),\ldots]$; ordinary zeta factors
in this coefficient ring are not counted as additional depth.

\begin{theorem}[Constructive depth bound]\label{spectral:depth}
Let $d\geq2$ be the largest number of factors in a monomial of $P$.
Formula \eqref{spectral:FP} has a finite expression in convergent multiple
zeta values and first derivatives in one coordinate of multiple-zeta germs,
all of depth at most $d-1$.  Every germ can be chosen holomorphic for
$\re s>-1$, with $s$ added to its first or second coordinate.
The principal coefficients and this finite conversion
are obtained by polynomial operations and finite stuffle products.
If $d=1$, the finite part in \eqref{spectral:FP} is zero.
\end{theorem}

\begin{proof}
Write $H_n^{(k)}=\sum_{a=1}^na^{-k}$, with $H_n=H_n^{(1)}$.
The exact local series is
\begin{equation}\label{spectral:localH}
 \psi(-n+v)=-\frac1v+H_n-\gamma
   +\sum_{k\geq1}
       \bigl(H_n^{(k+1)}+(-1)^{k+1}\zeta(k+1)\bigr)v^k.
\end{equation}
Differentiate and multiply to compute the finitely many required
$p_j(n)$.  Each term in a principal coefficient must use a singular
factor, so $p_j(n)$ is a polynomial of harmonic degree at most $d-1$.
Substitute $H_n^{(k)}=H_{n-1}^{(k)}+n^{-k}$ and subtract.  Thus
\begin{equation}\label{spectral:differencepoly}
 \Delta p_j(n)=\sum_{e\geq1}n^{-e}
             S_{j,e}\bigl(H_{n-1},H_{n-1}^{(2)},\ldots\bigr),
 \qquad \deg S_{j,e}\leq d-2,
\end{equation}
with only finitely many nonzero terms.  When $j\geq2$, the outer exponent
in \eqref{spectral:FP} is $e+j-1\geq2$.  Finite stuffle multiplication
therefore converts these terms to convergent multiple zeta values of
depth at most $d-1$.  The same argument applies to the terms $e\geq2$
in $D_P(s)$ and gives germs holomorphic for $\re s>-1$.

The possible $e=1$ terms require the residue cancellation, not an
unjustified separation into divergent multiple-zeta series.  Regard the
residue as the finite polynomial $R(X_1,X_2,\ldots)$ obtained from
\eqref{spectral:localH}.  Setting $X_k=\zeta(k)$ for $k\geq2$ turns that
local series into $X_1-\gamma+K(v)$.  Equation
\eqref{spectral:frozen} consequently gives the polynomial identity
\[
 R(X_1,\zeta(2),\zeta(3),\ldots)=0
 \quad\hbox{for every }X_1.
\]
No algebraic independence of zeta values is being assumed here.
The coefficient $S_{1,1}$ in \eqref{spectral:differencepoly} is
$\partial_{X_1}R$.  It belongs to the ideal generated by
$X_k-\zeta(k)$, $k\geq2$, and can be written as a sum of such factors
times polynomials of degree at most $d-3$; when $d=2$ this coefficient
vanishes.  On substitution each factor
has the convergent upper-tail expansion
\[
 H_{n-1}^{(k)}-\zeta(k)=-\sum_{a\geq n}a^{-k},\qquad k\geq2.
\]
Expand the remaining polynomial by finite stuffle products.  Each term
then uses a composition $\alpha$ of length at most $d-3$, and the exact
conversion is
\begin{equation}\label{spectral:tailconversion}
 \boxed{\displaystyle
 \sum_{n\geq1}\frac{(H_{n-1}^{(k)}-\zeta(k))H_{n-1}(\alpha)}{n^{1+s}}
     =-\zeta(k,1+s,\alpha)-\zeta(k+1+s,\alpha),\quad k\geq2.}
\end{equation}
Indeed, separate the upper-tail index $a\geq n$ into $a>n$ and $a=n$.
The largest resulting depth is $2+(d-3)=d-1$.  Every prefix sum of
real parts exceeds its prefix length when $\re s>-1$, so both series
on the right are absolutely convergent and holomorphic throughout that
half-plane.  The spectral parameter occupies the second coordinate in
the first series and the first coordinate in the second.  Differentiating
at $s=0$ completes the conversion of $D_P'(0)$.

If $d=1$, $P$ is a linear combination of positive derivatives of
digamma.  All its principal coefficients are constant in $n$, so
\eqref{spectral:FP} is zero.
\end{proof}

\begin{corollary}[Depth of the Gamma normal-form remainder]
\label{spectral:Gammadepth}
For $m\geq4$, write $\NN(\A_m)=u_0^m+R_m$.  The finite part of
$R_m(\psi,\psi',\ldots)$ has a constructive expression in the preceding
coordinates of depth at most $m-3$.
\end{corollary}

\begin{proof}
Every monomial of $R_m$ contains a positive derivative.  Its factor
degree has the same parity as $m$, and its weight is $m$; after removal
of $u_0^m$ the factor degree is therefore at most $m-2$.  Apply
Theorem~\ref{spectral:depth}.
\end{proof}

\subsection{Two explicit consistency checks}

For $P=u_1^2$, the nonzero principal coefficients are
\[
 p_4(n)=1,\qquad p_2(n)=2\bigl(H_n^{(2)}+\zeta(2)\bigr),\qquad
 r(n)=4\bigl(H_n^{(3)}-\zeta(3)\bigr).
\]
The general formula immediately yields
\begin{equation}\label{spectral:trigammasquare}
 \FP\int_0^1\psi'(x)^2\dd x
       =2\zeta(3)+4\zeta'(3).
\end{equation}
This is also a useful numerical check because all subtractions at the
origin are explicit.

For $P=u_0u_1^2$, the residue difference is
\[
 \Delta r(n)=\frac{4(H_{n-1}^{(3)}-\zeta(3))}{n}
       +\frac{2(H_{n-1}^{(2)}+\zeta(2))}{n^2}
       +\frac{4(H_{n-1}-\gamma)}{n^3}.
\]
The upper-tail conversion gives the fully convergent expression
\begin{align}\label{spectral:exampleDP}
 D_P(s)={}&-4\zeta(3,s+1)-4\zeta(s+4)
       +2\zeta(s+2,2)+2\zeta(2)\zeta(s+2)\notag\\
       &+4\zeta(s+3,1)-4\gamma\zeta(s+3),\qquad \re s>-1.
\end{align}
In particular,
\begin{equation}\label{spectral:digamma_trigammasquare}
 \FP\int_0^1\psi(x)\psi'(x)^2\dd x
       =\frac{35}{6}\zeta(4)-2\gamma\zeta(3)+D_P'(0).
\end{equation}
The term $\zeta(3,s+1)$ illustrates why the spectral derivative may
move to an interior coordinate.  Separately writing
$\zeta(s+1,3)-\zeta(3)\zeta(s+1)$ at $s=0$ would hide the cancellation
that the upper-tail argument makes explicit.

\section[An explicit sixth digamma moment and higher-weight consequences]{An explicit sixth digamma moment\\and higher-weight consequences}
\label{sixth:sec}

We now apply the two independent reductions: the scalar Gamma bridge
isolates the strict-depth Laurent slopes, while the differential normal
form leaves a small number of derivative-containing integrals.
At weight six only one genuinely mixed coordinate remains after the
quadratic polygamma formula has been used.  We specify that coordinate
completely, including ordinary convergent series and a finite
multiple-zeta representation.

\subsection{The mixed sixth-weight coordinate}

At $x=-n+v$ write
\[
 \psi(x)=-v^{-1}+b_n+c_nv+d_nv^2+e_nv^3+f_nv^4+O(v^5),
\]
where, only within this subsection,
\begin{equation}\label{sixth:locals}
 \begin{gathered}
 b_n=H_n-\gamma,\quad c_n=H_n^{(2)}+\zeta(2),\quad
 d_n=H_n^{(3)}-\zeta(3),\\
 e_n=H_n^{(4)}+\zeta(4),\quad f_n=H_n^{(5)}-\zeta(5).
 \end{gathered}
\end{equation}
These local $c_n$ are unrelated to the reciprocal-Gamma coefficients
$c_j$ of \eqref{gamma:c}.  Direct multiplication gives the principal
parts of $f=\psi^2\psi'^2$:
\begin{center}
\begin{tabular}{@{}cl@{}}
\toprule
$j$ & $p_j(n)$ in $f(-n+v)=\sum_{j=1}^6p_j(n)v^{-j}+O(1)$ \\
\midrule
6 & $1$ \\
5 & $-2b_n$ \\
4 & $b_n^2$ \\
3 & $-2b_nc_n+2d_n$ \\
2 & $2b_n^2c_n-2c_n^2-6b_nd_n+4e_n$ \\
1 & $4b_n^2d_n+2b_nc_n^2-6c_nd_n-10b_ne_n+6f_n$ \\
\bottomrule
\end{tabular}
\end{center}
Let $r_n=p_1(n)$ and define
\begin{equation}\label{sixth:D}
 \mathcal D(s)=\sum_{n\ge1}(r_n-r_{n-1})n^{-s},\qquad
 \Omega_5=\mathcal D'(0).
\end{equation}
By the residue theorem of Section~\ref{spectral:sec}, $\mathcal D$
is holomorphic for $\re s>-1$.  In particular these are ordinary
absolutely convergent identities:
\begin{equation}\label{sixth:omega-series}
 \boxed{\displaystyle
 \Omega_5=-\sum_{n\ge1}(r_n-r_{n-1})\log n
             =\sum_{n\ge1}r_n\log\left(1+\frac1n\right).}
\end{equation}
The last equality is summation by parts; its boundary term
$r_N\log(N+1)$ tends to zero.
The subscript five indicates the residual weight of the moment
calculation, not a Stieltjes index or a zeta derivative order.

For a short direct check, set $b=b_{n-1},c=c_{n-1},d=d_{n-1},e=e_{n-1}$.
Replacing the harmonic cutoff $n-1$ by $n$ in the residue polynomial
gives exactly
\begin{equation}\label{sixth:delta-r}
 r_n-r_{n-1}=
 \frac{8bd+2c^2-10e}{n}
 +\frac{4bc-2d}{n^2}
 +\frac{4b^2-2c}{n^3}-\frac4{n^5}.
\end{equation}
Its apparent $1/n$ term has a vanishing frozen coefficient.
Indeed, with $T_j=\zeta(j)-H_{n-1}^{(j)}$,
\begin{equation}\label{sixth:tail-cancel}
 8bd+2c^2-10e=-8bT_3-8\zeta(2)T_2+2T_2^2+10T_4,
\end{equation}
where $8\zeta(2)^2=20\zeta(4)$ was used.
This also directly proves the convergence asserted in
\eqref{sixth:D}.
Telescoping without logarithms gives the useful exact normalization
\begin{equation}\label{sixth:Dzero}
 \mathcal D(0)=4\gamma^2\zeta(3)-5\gamma\zeta(4)
                 -6\zeta(2)\zeta(3)+6\zeta(5).
\end{equation}

The finite multiple-zeta representation can be stated compactly as
$\mathcal D(s)=\sum_{\boldsymbol a}q_{\boldsymbol a}
\zeta(s+a_1,a_2,\ldots,a_d)$, with the following complete table.
The entries with first index one are combined before continuation
or differentiation; an individually convergent alternative is given
in Appendix~\ref{tables:omega-convergent}.

\begin{center}
\small
\begin{tabular}{@{}ll@{\qquad}ll@{}}
\toprule
$\boldsymbol a$ & $q_{\boldsymbol a}$ & $\boldsymbol a$ & $q_{\boldsymbol a}$\\
\midrule
$(1)$ & $8\gamma\zeta(3)+2\zeta(2)^2-10\zeta(4)$
 & $(2,2)$ & $-4\gamma$\\
$(1,1)$ & $-8\zeta(3)$ & $(2,2,1)$ & $4$\\
$(1,1,3)$ & $8$ & $(2,3)$ & $2$\\
$(1,2)$ & $4\zeta(2)$ & $(3)$ & $4\gamma^2-2\zeta(2)$\\
$(1,2,2)$ & $4$ & $(3,1)$ & $-8\gamma$\\
$(1,3)$ & $-8\gamma$ & $(3,1,1)$ & $8$\\
$(1,3,1)$ & $8$ & $(3,2)$ & $2$\\
$(2)$ & $-4\gamma\zeta(2)+2\zeta(3)$ & $(5)$ & $-4$\\
$(2,1)$ & $4\zeta(2)$ & $(2,1,2)$ & $4$\\
\bottomrule
\end{tabular}
\end{center}
To prove the table, expand \eqref{sixth:delta-r} and use only
$H_N^{(a)}H_N^{(b)}=H_N(a,b)+H_N(b,a)+H_N^{(a+b)}$.
No numerical recognition or conjectural relation is needed.

\begin{proposition}[The complete mixed integral]
\label{sixth:mixed}
With $\Omega_5$ defined by \eqref{sixth:D},
\begin{equation}\label{sixth:I}
 \boxed{\displaystyle
 \FP\int_0^1\psi(x)^2\psi'(x)^2\dd x
 =\Omega_5+2\gamma^2\zeta(3)-\frac{35}{3}\gamma\zeta(4)
      +\frac{31}{3}\zeta(2)\zeta(3)-\frac{19}{3}\zeta(5).}
\end{equation}
\end{proposition}
\begin{proof}
The nonresidue part of the general finite-part formula is
$J=\sum_{n\ge1}\sum_{j=2}^6
       (p_j(n)-p_j(n-1))/((j-1)n^{j-1})$.
The pole table reduces its entire summand to
\[
 j_n=\frac{4bc-6d}{n^2}+\frac{2b^2-3c}{n^3}
             -\frac{7b}{3n^4}-\frac{13}{6n^5}.
\]
Stuffle expansion expresses its sum in convergent multiple zeta values
of weight at most five.  The needed reductions are
\[
 \zeta(2,1)=\zeta(3),\quad
 \zeta(3,1)=\tfrac14\zeta(4),\quad
 \zeta(2,2)=\tfrac34\zeta(4),
\]
\begin{align*}
 \zeta(4,1)&=2\zeta(5)-\zeta(2)\zeta(3),\\
 \zeta(3,2)&=3\zeta(2)\zeta(3)-\tfrac{11}{2}\zeta(5),\\
 \zeta(2,3)&=\tfrac92\zeta(5)-2\zeta(2)\zeta(3),
\end{align*}
together with the dualities
$\zeta(3,1,1)=\zeta(4,1)$,
$\zeta(2,2,1)=\zeta(3,2)$, and
$\zeta(2,1,2)=\zeta(2,3)$.
These classical low-weight identities can be obtained by the
iterated-integral substitution $t\mapsto1-t$ and the shuffle/stuffle
product rules.  For example the weight-five double sum formula,
the stuffle of $\zeta(2)\zeta(3)$, and its shuffle give respectively
\[
 \zeta(2,3)+\zeta(3,2)+\zeta(4,1)=\zeta(5),
\]
\[
 \zeta(2)\zeta(3)=\zeta(2,3)+\zeta(3,2)+\zeta(5)
                =\zeta(2,3)+3\zeta(3,2)+6\zeta(4,1),
\]
which are a nonsingular linear system for the displayed evaluations.
Substitution in $\sum j_n$ gives
$J=2\gamma^2\zeta(3)-35\gamma\zeta(4)/3
   +31\zeta(2)\zeta(3)/3-19\zeta(5)/3$.
The residue contribution is $\mathcal D'(0)=\Omega_5$, proving
\eqref{sixth:I}.
\end{proof}

\subsection{The sixth power with no unnamed mixed integral}

Return to $c_j=[u^j]\Gamma(1+u)^{-1}$.  The required coefficients are
\begin{align}
 c_2&=\frac{\gamma^2-\zeta(2)}2,\notag\\
 c_3&=\frac{\gamma^3-3\gamma\zeta(2)+2\zeta(3)}6,\label{sixth:cs}\\
 c_4&=\frac{\gamma^4-6\gamma^2\zeta(2)+3\zeta(2)^2
                     +8\gamma\zeta(3)-6\zeta(4)}{24}.\notag
\end{align}

\begin{theorem}[The sixth digamma moment]
\label{sixth:main}
In the original coordinate finite-part convention,
\begin{equation}\label{sixth:Q6}
 \boxed{\begin{aligned}
 \FP\int_0^1\psi(x)^6\dd x={}&
 15\Omega_5-720\lambda_5+720\gamma\lambda_4
 -720c_2\lambda_3+720c_3\lambda_2+720c_4\gamma_1\\
 &-60\gamma^3\zeta(2)+180\gamma^2\zeta(3)-70\gamma\zeta(4)\\
 &-55\zeta(2)\zeta(3)+175\zeta(5)
                  +48\zeta'(5)-135\zeta(6).
 \end{aligned}}
\end{equation}
All constants on the right have specified ordinary-series, Laurent,
or standard special-function definitions.  In particular, this
formula makes no assertion that $\Omega_5$ reduces to a smaller
constant alphabet.
\end{theorem}
\begin{proof}
The exact primitive in the sixth-weight normal decomposition is
\begin{equation}\label{sixth:H6}
 H_6=-3u_0^5+20u_0^3u_1-15u_0^2u_2-15u_0u_1^2
                       +6u_0u_3+9u_1u_2-u_4.
\end{equation}
Direct differentiation gives
$\A_6=u_0^6-15u_0^2u_1^2+u_2^2+DH_6$.
The Laurent expansion \eqref{gamma:local} gives
\begin{equation}\label{sixth:boundary}
 \Delta H_6=60\gamma^3\zeta(2)-150\gamma^2\zeta(3)
 -105\gamma\zeta(4)+210\zeta(2)\zeta(3)-242\zeta(5).
\end{equation}
The scalar bridge is
\[
 \mathcal L_6=-135\zeta(6)
  -720\{c_4\lambda_1-c_3\lambda_2+c_2\lambda_3
                              -\gamma\lambda_4+\lambda_5\}.
\]
The Hurwitz-product formula gives
$\FP\int\psi''(x)^2\dd x=-28\zeta(5)-48\zeta'(5)$.
Therefore
\[
 Q_6=\mathcal L_6+15\FP\int_0^1\psi^2\psi'^2\dd x
                    -\FP\int_0^1\psi''^2\dd x-\Delta H_6.
\]
Substitute \eqref{sixth:I}, \eqref{sixth:boundary}, and
$\lambda_1=-\gamma_1$; collecting rational coefficients proves
\eqref{sixth:Q6}.
\end{proof}

As numerical reference values, independently evaluated quantities are
\begin{align*}
 \Omega_5&=-10.13986531321045747401304000835176665\ldots,\\
 \FP\int_0^1\psi^2\psi'^2\dd x
   &=-2.76249428761727870446208254860961992\ldots,\\
 Q_6&=-64.95511203538334575217535095818579002\ldots.
\end{align*}
The mixed-identity residual was below $2\cdot10^{-99}$ in the
115-digit run.  An end-to-end check of \eqref{sixth:Q6}, using
independent Mellin integrals for all five $\lambda_d$, had residual
below $10^{-87}$ at 90 working digits.  These are reproducible
floating-point diagnostics, not certified error enclosures.

\subsection{Seventh, eighth, and arbitrary weights}

The next two normal forms in \eqref{gamma:firstforms} give
\begin{align}
 Q_7={}&\mathcal L_7-\Delta H_7
    +35\FP\int_0^1\psi^3\psi'^2\dd x
    -7\FP\int_0^1\psi\psi''^2\dd x,\label{sixth:Q7}\\
 Q_8={}&\mathcal L_8-\Delta H_8
    +70\FP\int_0^1\psi^4\psi'^2\dd x
    -28\FP\int_0^1\psi^2\psi''^2\dd x\notag\\
 &\hspace{14mm}+7\FP\int_0^1\psi'^4\dd x
    +\FP\int_0^1\psi'''^2\dd x.\label{sixth:Q8}
\end{align}
These are useful compact reductions rather than claims of minimal
arithmetic bases.  Every remaining integral is evaluated by the
general pole-coefficient procedure, and the last square is already
the explicit quadratic Hurwitz value.
The endpoint polynomials $\Delta H_7,\Delta H_8$ appear in
Appendix~\ref{tables:boundaries}; the data file retains all exact
coefficients without a choice of even-zeta simplification.

At general weight $m$, the closed coefficient sectors determine every
term $\psi^{m-2r-2}(\psi^{(r)})^2$ and every term containing only
$\psi$ and $\psi'$.  The remaining functional monomials can be
generated finitely.  Their residue differences have finite harmonic
word expansions, and their antiderivatives follow from the normally
convergent primitive theorem.  This supplies a systematic way to
produce exact identities and isolate the next arithmetic question at
each weight.  It does not turn formal dimension into a theorem about
independence of numerical constants.

\section{Source audit, proposed corrections, and further research}
\label{audit:sec}

\subsection{What the audit establishes}

The repository was inspected at the revision on the title page.
All six incoming ZIPs were unpacked with a path-safety check; their
inventories, claim/status files, and relevant sections were read.
The accompanying \texttt{incoming\_audit\_summary.json} records the
archive names, SHA-256 digests, already covered results, and open
questions.  This was a targeted mathematical audit of the claims
used here, not a line-by-line certification of every theorem in the
large source corpus.

The source claims used most directly are the exact even-tail
Bernoulli coordinates, the first-$q$ conjecture, the Gamma-quotient
ordinary integral, the fifth-moment baseline, and the recorded
Family--2 example.  The tail coefficient/operator conversion was
checked independently; the new proof includes the column reversal
and the omitted lower-index-zero convention.  The Gamma bridge
was reproved with its Laurent normalization before it was used in
the sixth moment.  The Family--2 equation was compared to its
original right-hand side, not merely to a numerically similar
simplification.

\textbf{No new false theorem in the inherited material was established
by this audit.}  Some incoming reports already correct signs,
Stieltjes indexing, or expressibility claims in external papers.
Those corrections are identified as inherited in the audit record;
they are not counted as new findings here.

\subsection{Concrete proposed source changes}

\begin{enumerate}
\item \textbf{Upgrade the Family--2 status.}
In the canonical \texttt{chapters/04-depth.tex}, the equation labeled
\texttt{double:eq:fam2-example} can be marked proved.
Retain its original identity for continuity, add
\eqref{bip:eq:family2-short} as its simplified form, and cite the
two exact component evaluations \eqref{bip:eq:12}--\eqref{bip:eq:21}.
The discovery by integer relation remains a useful historical note,
but it is no longer the justification for that equation.

\item \textbf{Refine the optimal-prefix conjecture.}
In the incoming canonical fifth-moment report, retain
\texttt{conj:firstq} for general unequal even orders.
Record $p=q-2$ as a theorem with determinant
\eqref{adj:determinant}; retain $p=q$ as the already proved source
case and general gaps $q-p\ge4$ as open.
This is an infinite-family resolution, not a complete proof of the
original conjecture.

\item \textbf{Attach the sixth calculation to the higher-moment agenda.}
The new result \eqref{sixth:Q6} completes a specific sixth-power
calculation in explicit convergent zeta-derivative coordinates.
Its remaining arithmetic question is the reduction of $\Omega_5$.
The generic all-power algorithm and the fifth calculation were
already present and should retain their earlier attribution.

\item \textbf{Cite the exact antecedents for two reusable methods.}
The functional-monomial direct sum should cite
\cite[Proposition~4.10]{Guo2014}; the bicolored parity implementation
should cite Nakamura and Panzer \cite{bip:Nakamura2012,bip:Panzer2017}.
The latter paper's footnote~16 identifies a transposed product in
Nakamura's printed equation (1.3).  Our proof does not import that
formula without rederivation.  This is an attribution and formula-use
clarification, not an assertion that the supplied manuscript copied
the misprint.

\item \textbf{Preserve the precise meanings of ``finite part'' and ``rank''.}
When importing the new results, keep the original-coordinate
definition \eqref{scope:finitepart}, the local correction in the
Hurwitz product section, and the formal coefficient-space qualifiers.
A raw Laurent constant of the spectral kernel generally differs
from the coordinate integral; $\NN$ is not multiplicative; and a
formal rank computation does not prove arithmetic independence.
These are explicit integration requirements for this article, not
newly alleged errors in every prior use of those terms.
\end{enumerate}

\subsection{Research questions with concrete mathematical targets}

The following agenda separates exact next calculations from genuine
structural questions.  An answer should supply a proof, a fully stated
functional identity, or an exact algebraic certificate.  A large
precision search or a numerical rank alone would not settle the
corresponding question.

\paragraph{1. The first unresolved even-tail gap.}
Prove or disprove the optimal first-$q$ conjecture for
$q=p+4$, uniformly in $p$.  The adjacent-order proof uses one free
formal coefficient and a sign-definite scalar recurrence.  The next
gap produces a larger initial-data space.  The concrete target is a
nonzero determinant formula for the associated finite block, perhaps
through a matrix recurrence whose relevant minors have a controlled
sign.  This would test whether the adjacent argument extends or
depends essentially on second order.

\paragraph{2. A determinant formula for arbitrary unequal orders.}
For $p=2a,q=2b$, seek a closed factorization or a finite positive
coefficient formula for the late Bernoulli matrix determinant.
Any proposed expression must retain the source's omitted binomial
terms and the centered cutoff.  A proof should also give the inverse
map, thereby producing ordinary harmonic-tail identities in addition
to a rank statement.

\paragraph{3. Arithmetic structure of the adjacent determinant sequence.}
The rational numbers $c_n$ have both a Bernoulli recurrence and the
central-binomial representation \eqref{adj:detgenerator}.
Determine an exact denominator formula, a useful integer
renormalization, or a finite convolution identity for this sequence.
The target is an arithmetic identity or recurrence with transparent
integrality, rather than an asymptotic estimate.

\paragraph{4. The next full Gamma coefficient sector.}
Find the generating function for all functional monomials containing
$u_2^2$ and arbitrary powers of $u_0,u_1$, including any terms
coupled to $u_2^3$ by factor-degree parity.
The ideal argument used for the complete $u_0,u_1$ sector suggests
working modulo a precisely stated higher-derivative monomial ideal.
Its preservation by the normal projection must be proved from the
reduction rule; stability under $D$ by itself is insufficient.

\paragraph{5. Higher sectors and a global normal-form generating function.}
The centered generator contains only even derivative indices before
normal reduction.  Can that symmetry be turned into a closed
multivariable generating function for $\NN(\A_m)$, beyond the
square and digamma--trigamma sectors?  An exact coefficient-extraction
formula indexed by partitions would be useful even without a short
elementary expression.

\paragraph{6. Arithmetic reduction of the mixed sixth coordinate.}
Determine whether $\Omega_5$ can be expressed in a prescribed
algebra generated by the height-one $\lambda_d$, ordinary zeta
derivatives, and lower mixed coordinates already appearing in the
fifth moment.  Appendix~\ref{tables:omega-convergent} supplies a
finite list of absolutely convergent MZV germs of depth at most three.
This makes the question precise enough for a symbolic functional
relation search.  Failure of a finite search would not prove
nonmembership or independence.

\paragraph{7. Complex-order functional relations before differentiation.}
The preceding question is about derivatives in zeta orders, not only
values at integral orders.  A parity identity known only for integer
indices cannot be differentiated in an index without an analytic
extension.  Seek functional equations for the required complex-order
germs, prove their continuation and local regularity, and only then
differentiate.  This is a specific way to avoid an unjustified
``differentiate the MZV reduction'' step.

\paragraph{8. Minimal spectral depth of higher moments.}
The normal-form and residue theorems give an upper bound on the
depth needed for the derivative-containing remainder of $Q_m$.
Determine when that bound is attained after all proven functional
relations are used.  First complete explicit seventh- and eighth-weight
coordinate tables, then look for uniform cancellations.
Distinguish the dimension of differential-polynomial normal forms,
the rank of their harmonic residue maps, and the dimension of their
actual period values; these are different problems.

\paragraph{9. Generalized Stieltjes products at the missing derivative edge.}
The finite generator in \eqref{hprod:stieltjes-family} treats every
pair of positive argument-derivative orders.  Extend its manifestly
holomorphic coefficient form to cases with one or both derivative
orders zero, and prove compatibility with the incoming
Stieltjes-collision formulas.  The extra spectral poles at order one
must be subtracted jointly; blindly setting an integer parameter to
one in the positive-order formula does not establish this extension.

\paragraph{10. Exact cyclotomic identities in the complementary parity sector.}
The bicolored theorem computes real parts in odd weight and imaginary
parts in even weight.  The Gaussian $S_6$ and revised $S_8$ questions
are not resolved by that mechanism.  Seek a genuinely additional
relation: a proven depth-three reduction, an integral symmetry, or a
specialization of a functional equation with a larger argument set.
The precise definitions and corrected candidates in the current
source should be retained; already rejected coefficient lists should
not be recycled as conjectures.

\paragraph{11. Antiderivatives and a specified function class.}
The spectral theorem provides normally convergent anchored primitives
for every derivative-containing differential polynomial.
Which of these primitives reduce to a specified finite class of
Gamma, Barnes, polylogarithm, and zeta-order-derivative functions?
A useful first target is the primitive of $\psi^2\psi'^2$.
Any nonexistence claim must name the allowed class and prove an
obstruction within it; an unspecified notion of ``closed form''
does not define a mathematical problem.

\paragraph{12. Exact proof certificates for repository integration.}
Turn the polynomial reductions, low-weight shuffle identities, and
finite Bernoulli matrix conjugacies into machine-readable exact
certificates.  The current scripts already separate rational
identities from numerical diagnostics.  A next stage could check
those certificates with an independent implementation or a proof
assistant, while leaving analytic continuation and contour estimates
as explicit lemmas with their assumptions stated.

The most immediate next proof target is the $q-p=4$ tail determinant.
The most immediate next arithmetic target is the convergent
depth-three expression for $\Omega_5$.  Both have precise finite
input data in this article, and both would produce exact identities
of direct use to the existing manuscript.

\appendix
\section{Exact endpoint and convergent-coordinate tables}
\label{tables:sec}

\subsection{Endpoint corrections through weight eight}
\label{tables:boundaries}

For the unique primitives $H_m=\HH(\A_m)$ in
\eqref{gamma:decomp}, the endpoint functional is
\begin{align*}
 \Delta H_2={}&0,\\
 \Delta H_3={}&-3\zeta(2),\\
 \Delta H_4={}&12\gamma\zeta(2)-10\zeta(3),\\
 \Delta H_5={}&-30\gamma^2\zeta(2)+50\gamma\zeta(3)
                                      +\tfrac{35}{2}\zeta(4),\\
 \Delta H_6={}&60\gamma^3\zeta(2)-150\gamma^2\zeta(3)
 -105\gamma\zeta(4)+210\zeta(2)\zeta(3)-242\zeta(5),\\
 \Delta H_7={}&-105\gamma^4\zeta(2)+350\gamma^3\zeta(3)
                     +\tfrac{735}{2}\gamma^2\zeta(4)\\
 &-1470\gamma\zeta(2)\zeta(3)+1694\gamma\zeta(5)
                +469\zeta(3)^2-\tfrac{1085}{2}\zeta(6),\\
 \Delta H_8={}&168\gamma^5\zeta(2)-700\gamma^4\zeta(3)
                          -980\gamma^3\zeta(4)\\
 &+5880\gamma^2\zeta(2)\zeta(3)-6776\gamma^2\zeta(5)
             -3752\gamma\zeta(3)^2+4340\gamma\zeta(6)\\
 &+7840\zeta(2)\zeta(5)-1862\zeta(3)\zeta(4)-11824\zeta(7).
\end{align*}
They use only the elementary simplifications
$\zeta(2)^2=5\zeta(4)/2$,
$\zeta(2)^3=35\zeta(6)/8$, and
$\zeta(2)\zeta(4)=7\zeta(6)/4$.
The data file \texttt{inverse\_gamma\_normal\_forms.json} retains
the original independent-symbol expressions as well as each exact
primitive and normal form.  Repeating the finite Laurent coefficient
extraction proves every entry of this table.

\subsection{A fully convergent 17-term representation of \texorpdfstring{$\mathcal D(s)$}{D(s)}}
\label{tables:omega-convergent}

The first-index-one table in Section~\ref{sixth:sec} is convenient for
formal harmonic-word generation, but individual entries there may
have spectral poles.  The following alternative avoids them entirely.
The function $\mathcal D(s)$ is the sum of each function in the left
column multiplied by its coefficient in the right column:
\begin{center}
\begin{tabular}{@{}ll@{}}
\toprule
Individually convergent MZV germ & Coefficient \\
\midrule
$\zeta(2+s)$ & $2\zeta(3)-4\gamma\zeta(2)$\\
$\zeta(3+s)$ & $4\gamma^2-8\zeta(2)$\\
$\zeta(4+s)$ & $8\gamma$\\
$\zeta(5+s)$ & $6$\\
$\zeta(2+s,1)$ & $4\zeta(2)$\\
$\zeta(2+s,2)$ & $-4\gamma$\\
$\zeta(2+s,3)$ & $2$\\
$\zeta(3+s,1)$ & $-8\gamma$\\
$\zeta(4+s,1)$ & $-8$\\
$\zeta(2+s,1,2)$ & $4$\\
$\zeta(2+s,2,1)$ & $4$\\
$\zeta(3+s,1,1)$ & $8$\\
$\zeta(2,1+s)$ & $-6\zeta(2)$\\
$\zeta(3,1+s)$ & $8\gamma$\\
$\zeta(4,1+s)$ & $10$\\
$\zeta(2,1+s,2)$ & $-2$\\
$\zeta(3,1+s,1)$ & $-8$\\
\bottomrule
\end{tabular}
\end{center}
Every entry converges absolutely and defines a holomorphic function
on $\re s>-1$.  Therefore $\Omega_5$ is also the sum of the
\emph{individual ordinary derivatives at zero} of these 17 germs,
with the same coefficients.  Derivatives in the second coordinate
must remain second-coordinate derivatives.

For an exact derivation, use the elementary identity, valid for
$k\ge2$ and a positive-index word $\boldsymbol a$,
\begin{equation}\label{tables:upper-tail}
 \begin{split}
 &\sum_{n\ge1}\frac{(H_{n-1}^{(k)}-\zeta(k))
                             H_{n-1}(\boldsymbol a)}{n^{1+s}}\\
 &\hspace{18mm}=-\zeta(k,1+s,\boldsymbol a)
                -\zeta(k+1+s,\boldsymbol a).
 \end{split}
\end{equation}
It follows by writing the negative tail as
$-\sum_{j\ge n}j^{-k}$ and partitioning $j>n$ from $j=n$.
All terms are absolutely summable in the stated half-plane; any
additional inner unit indices contribute at most logarithms.
The coefficient of $1/n$ in \eqref{sixth:delta-r} factors exactly as
\[
 8(H_{n-1}-\gamma)(H_{n-1}^{(3)}-\zeta(3))
 +2(H_{n-1}^{(2)}-\zeta(2))(H_{n-1}^{(2)}+3\zeta(2))
 -10(H_{n-1}^{(4)}-\zeta(4)).
\]
Apply \eqref{tables:upper-tail} to these three terms and retain
the first-index-at-least-two rows of the original table.
Two $\zeta(3+s,2)$ terms cancel; collecting the single zeta terms
gives precisely the 17 rows above.  This establishes the alternative
without analytic regularization of any individual summand.

The same residue calculation gives an exact additional continuation
statement.  Its first two asymptotic coefficients are
$r_n-r_{n-1}=8\zeta(2)(\log n-1)/n^2+
O(\log^2(n+2)/n^3)$, obtained by substituting the ordinary harmonic
expansions in \eqref{sixth:delta-r}.  Consequently
\[
 \mathcal D(s)+8\zeta(2)\{\zeta'(s+2)+\zeta(s+2)\}
\]
is holomorphic for $\re s>-2$.  In particular the exact principal
part of $\mathcal D$ at $s=-1$ is
\[
 \frac{8\zeta(2)}{(s+1)^2}-\frac{8\zeta(2)}{s+1}.
\]
The asymptotic estimate here is used only to prove that spectral
identity and its pole coefficients.

\subsection{A convergent logarithmic antiderivative for the mixed term}

For $a>0$ and $\re z>0$, the pole table in Section~\ref{sixth:sec}
gives the anchored primitive
\begin{equation}\label{tables:primitive6}
 \begin{split}
 \mathcal P_6(z;a)=\sum_{n\ge0}\biggl\{
 &r_n\bigl(\Log(n+z)-\log(n+a)\bigr)\\
 &+\sum_{j=2}^6\frac{p_j(n)}{1-j}
       \bigl((n+z)^{1-j}-(n+a)^{1-j}\bigr)\biggr\}.
 \end{split}
\end{equation}
It converges normally on compact subsets of the right half-plane,
vanishes at $z=a$, and satisfies
$\partial_z\mathcal P_6(z;a)=\psi(z)^2\psi'(z)^2$.
Thus the mixed sixth coordinate has a fully specified antiderivative
as well as its endpoint finite-part value.  The general primitive
theorem supplies the same construction for all the derivative terms
in the seventh, eighth, and higher moment reductions.

\section{Reproducibility and the limits of computation}
\label{repro:sec}

The package separates proof, exact verification, and floating-point
diagnostics.  The proofs of the infinite families appear in the
article.  Exact checks compare independent finite constructions or
verify finite polynomial identities.  Numerical residuals test the
analytic formulas at selected instances; they are not interval
certificates and do not prove arbitrary-order assertions.

\begin{center}
\small
\begin{tabularx}{\textwidth}{@{}p{.25\textwidth}X@{}}
\toprule
Verification family & Completed evidence \\
\midrule
Bicolored parity & 49 exact partial-fraction checks; 48 independent
Mellin-integral parity comparisons, plus explicit specializations and
the original Family--2 right-hand side.  The largest parity residual
was $5.94\cdot10^{-50}$ at 50 working digits. \\
Adjacent tails & 127 exact assertions, including matrix/operator
conjugacy and determinants through $b=16$, recurrence signs through
$c_{40}$, and source-coordinate cancellation.  Eight numerical
diagnostics had maximum scaled residual $1.62\cdot10^{-77}$. \\
Gamma sectors & Independent Bell-polynomial and normal-form construction
through weight 16; 49 square coefficients, 17 full sectors, and the
hyperbolic-sine-integral generator agree exactly. \\
Hurwitz products & Four independent local-subtraction quadratures and
one logarithmic order-derivative quadrature, all at 85 working digits;
absolute residuals were below $7\cdot10^{-85}$. \\
Spectral polynomials & Independent principal-part checks, frozen-residue
cancellation, the exact mixed residue difference, and a separately
regularized trigamma-square quadrature. \\
Mixed sixth coordinate & Six principal parts, a Laurent constant,
two formal harmonic-word tables, 24 finite-harmonic checks, 54 finite
tail-ordering checks, and the exact weight-five reduction.  The
115-digit mixed identity residual was $1.76\cdot10^{-99}$. \\
Full sixth moment & Independent local-series quadrature for $Q_6$ and
Mellin integrals for all five strict-depth slopes.  The 90-digit
end-to-end residual was $9.90\cdot10^{-88}$. \\
\bottomrule
\end{tabularx}
\end{center}

The scripts require Python 3, SymPy, and mpmath.  The actual run used
SymPy 1.14.0 and mpmath 1.3.0.  Exact arithmetic uses rational numbers
or independent symbolic constants.  In the numerical tail calculations,
the finite sum and the asymptotic Hurwitz tail are evaluated with
separately specified cutoffs and working precision; the JSON records
retain those parameters and the observed discrepancies.

The package includes:
\begin{itemize}
\item the compiled PDF, one self-contained TeX source, and modular
TeX sources with stable theorem labels;
\item exact Gamma normal-form and mixed-coordinate data;
\item independent verification scripts and their completed JSON records;
\item the pinned source/archive audit, a claim ledger, and proposed
integration changes;
\item a build command, dependency notes, and a package checksum manifest.
\end{itemize}

No original incoming report is silently replaced.  No source conjecture
outside the explicitly proved cases is marked solved.  A suggested
integration is to add the article as a new incoming research report,
then promote the two proved source claims using the precise statements
in Section~\ref{audit:sec}.  The source revision and archive digests
make that comparison reproducible even if the repository subsequently
changes.

\begin{thebibliography}{99}
\raggedright

\bibitem{SourceManuscript}
ProveIt project,
\emph{Polylogarithms manuscript}, canonical modular manuscript,
\texttt{Analysis/Polylogarithms/docs/manuscript},
repository revision
\texttt{f7e35be357654198c1546fd40f363694c5bcc2cc}.
\url{https://github.com/VladimirReshetnikov/ProveIt}.
The source equation \texttt{double:eq:fam2-example} is in
\texttt{chapters/04-depth.tex}; the coincident Hurwitz kernel is in
\texttt{chapters/07-shifted-stieltjes-correlations.tex}.

\bibitem{SourceCanonical}
ProveIt project,
\emph{Canonical Digamma Polynomials and the Fifth Harmonic Bridge},
incoming report in
\path{ProveIt_Canonical_Digamma_Fifth_Moment_2026-10-11.zip},
at the same pinned revision.
Sections 3, 7, and 9 contain the Gamma bridge, exact even-tail
coordinates, and Conjecture \texttt{conj:firstq}.
Archive SHA-256:
\path{3c2fa00434f0366089b819bad77c5a176443594ed1fe6dd0731c877bc74d01a7}.

\bibitem{SourceHigher}
ProveIt project,
\emph{Higher Digamma Powers and Ordered Gamma Sums},
incoming report in
\path{ProveIt_Higher_Digamma_Powers_2026-10-11.zip},
at the same pinned revision.
The all-power algorithm, fifth moment, and normalized primitives are
prior baseline results.

\bibitem{bip:Nakamura2012}
T. Nakamura,
\emph{A simple proof of the functional relation for the Lerch type
Tornheim double zeta function},
Tokyo Journal of Mathematics \textbf{35} (2012), 333--337.
\url{https://arxiv.org/abs/1012.1144}.
See Proposition 1.2, read with the product correction recorded by
Panzer in footnote 16 of the next reference.

\bibitem{bip:Panzer2017}
E. Panzer,
\emph{The parity theorem for multiple polylogarithms},
Journal of Number Theory \textbf{172} (2017), 93--113.
\url{https://arxiv.org/abs/1512.04482}.
See Theorem 1.3, equation (3.2), and footnote 16.

\bibitem{Guo2014}
L. Guo, G. Regensburger, and M. Rosenkranz,
\emph{On integro-differential algebras},
Journal of Pure and Applied Algebra \textbf{218} (2014), 456--473.
\url{https://doi.org/10.1016/j.jpaa.2013.06.015};
\url{https://arxiv.org/abs/1212.0266}.
Proposition 4.10 gives the functional-monomial direct-sum decomposition.

\bibitem{Rosenkranz2008}
M. Rosenkranz and G. Regensburger,
\emph{Integro-differential polynomials and operators},
Proceedings of ISSAC 2008, 261--268.
The reduction rule is given in Lemma 5; the later treatment
\cite{Guo2014} supplies the direct-sum formulation used here.

\bibitem{Shpot2016}
M. A. Shpot, M. P. Chaudhary, and R. B. Paris,
\emph{Integrals of products of Hurwitz zeta functions and the Casimir
effect in $\phi^4$ field theories},
arXiv:1603.00722v2 (2016).
\url{https://arxiv.org/abs/1603.00722}.
Remark 1 states the classical Mikol\'as product integral.  Theorem 2
concerns the shifted, ordinary product integral, which is distinct
from the coordinate finite part evaluated in this article.

\bibitem{DLMF}
NIST Digital Library of Mathematical Functions,
\url{https://dlmf.nist.gov/}.
Relevant formulas and sections: 5.5.4 (digamma reflection),
5.7 (Gamma and digamma series), 5.11.2 and 5.15.9
(digamma and polygamma asymptotics), 24.8.1
(Bernoulli Fourier series), and 25.11 (Hurwitz zeta).
Accessed 11 October 2026.

\end{thebibliography}

\end{document}
